\documentclass[12pt]{article}
\linespread{1.25}
\pagestyle{plain}
\usepackage{graphicx}
\usepackage[export]{adjustbox}
\graphicspath{ {./images/} }
\title{{\Huge{\textbf{Advanced Financial Models}}}}
\author{Wei Chen}
\date{\today}
% -------------------------------------------------
% Fonts, language, and general layout
% -------------------------------------------------
\usepackage[T1]{fontenc}
\usepackage[english]{babel}
\usepackage{mathpazo}
\linespread{1.5}

\usepackage[a4paper,
  top=0.5in,
  bottom=0.2in,
  left=0.5in,
  right=0.5in,
  footskip=0.3in,
  includefoot
]{geometry}

\pagestyle{plain}
\setlength{\parindent}{0.2in}
\setlength{\parskip}{0pt}
\setlength{\columnseprule}{0pt}

% -------------------------------------------------
% Core packages
% -------------------------------------------------
\usepackage{graphicx}
\usepackage[export]{adjustbox}
\usepackage{array}
\usepackage{wrapfig}
\usepackage{multicol}
\usepackage{framed}
\usepackage{indentfirst}
\usepackage[normalem]{ulem}
\usepackage[inline,shortlabels]{enumitem}
\usepackage[dvipsnames]{xcolor}
\usepackage[bottom,multiple]{footmisc}

% -------------------------------------------------
% Mathematics
% -------------------------------------------------
\usepackage{amsmath,amsthm,amssymb}
\usepackage{mathtools}
\usepackage{mathrsfs}

% -------------------------------------------------
% TikZ and graphics
% -------------------------------------------------
\usepackage[most]{tcolorbox}
\usepackage{tikz}
\usepackage{tikz-3dplot}
\usepackage{tikz-cd}
\usepackage{tkz-tab}
\usepackage{tkz-euclide}
\usepackage{pgf}
\usepackage{pgfplots}
\pgfplotsset{compat=newest}

% -------------------------------------------------
% Hyperlinks and references
% -------------------------------------------------
\usepackage{hyperref}
\hypersetup{
  hypertex=true,
  colorlinks=true,
  linkcolor=blue,
  anchorcolor=blue,
  citecolor=blue
}
\usepackage[nameinlink]{cleveref}

% -------------------------------------------------
% Section and appendix formatting
% -------------------------------------------------
\usepackage[title,titletoc]{appendix}
\usepackage[explicit]{titlesec}

\titleformat{\chapter}
  {\fontsize{22}{26}\normalfont\bfseries}
  {Chapter \thechapter}
  {22pt}
  {\uline{#1}}

\titleformat{\section}
  {\fontsize{18}{30}\selectfont\bfseries}
  {\thesection}
  {16pt}
  {#1}

\titleformat{\subsection}
  {\fontsize{14}{18}\sffamily\bfseries}
  {\thesubsection}
  {12pt}
  {#1}

\titleformat{\subsubsection}
  {\fontsize{10}{12}\sffamily\large\bfseries}
  {\thesubsubsection}
  {8pt}
  {#1}

\titlespacing*{\section}{0pt}{5pt}{5pt}
\titlespacing*{\subsection}{0pt}{5pt}{5pt}
\titlespacing*{\subsubsection}{0pt}{5pt}{5pt}

% -------------------------------------------------
% Lists
% -------------------------------------------------
\setlist[enumerate,1]{label=(\alph*)., ref=\alph*}
\setlist{topsep=2pt,itemsep=2pt,parsep=0pt,partopsep=0pt}

% -------------------------------------------------
% Common shortcuts
% -------------------------------------------------
\newcommand{\remind}[1]{\textcolor{red}{\textbf{#1}}}
\newcommand{\hide}[1]{}
\newcommand{\Disp}{\displaystyle}
\newcommand{\qe}{\hfill\(\bigtriangledown\)}
\newcommand{\mt}[1]{\texorpdfstring{$#1$}{Text}}
\newcommand{\wt}[1]{\widetilde{#1}}
\newcommand{\textblue}[1]{\textbf{\textcolor{blue}{#1}}}

% -------------------------------------------------
% Paired delimiters and notation
% -------------------------------------------------
\newcommand{\norm}[1]{\left\lVert#1\right\rVert}
\newcommand{\inner}[2]{\left\langle #1, #2 \right\rangle}
\DeclarePairedDelimiter{\paren}{(}{)}
\DeclarePairedDelimiter{\brackets}{[}{]}
\DeclarePairedDelimiter{\braces}{\lbrace}{\rbrace}
\DeclarePairedDelimiter{\abs}{\lvert}{\rvert}

\newcommand{\floor}[1]{\left\lfloor #1 \right\rfloor}
\newcommand{\ceil}[1]{\left\lceil #1 \right\rceil}
\newcommand{\limit}[2]{\lim\limits_{#1 \rightarrow #2}}

% -------------------------------------------------
% Blackboard bold and calligraphic symbols
% -------------------------------------------------
\newcommand{\R}{\mathbb{R}}
\newcommand{\N}{\mathbb{N}}
\newcommand{\Q}{\mathbb{Q}}
\newcommand{\C}{\mathbb{C}}
\newcommand{\D}{\mathbb{D}}
\newcommand{\E}{\mathbb{E}}
\newcommand{\pr}{\mathbb{P}}
\newcommand{\eh}{\mathbb{H}}

\newcommand{\X}{\mathcal{X}}
\newcommand{\F}{\mathcal{F}}
\newcommand{\A}{\mathcal{A}}
\newcommand{\B}{\mathcal{B}}
\newcommand{\M}{\mathcal{M}}
\newcommand{\G}{\mathcal{G}}
\newcommand{\es}{\mathcal{S}}
\newcommand{\ml}{\mathcal{L}}

% -------------------------------------------------
% Text and symbol shortcuts
% -------------------------------------------------
\newcommand{\fat}{\operatorname{fat}}
\newcommand{\VC}{\text{VC}}
\newcommand{\var}{\operatorname{Var}}
\newcommand{\prob}{\operatorname{Prob}}
\newcommand{\argmin}{\operatornamewithlimits{argmin}}
\newcommand{\argmax}{\operatornamewithlimits{argmax}}

\newcommand{\ep}{\varepsilon}
\newcommand{\vp}{\varphi}
\newcommand{\lam}{\lambda}
\newcommand{\Lam}{\Lambda}
\newcommand{\ka}{\kappa}
\newcommand{\e}{\varnothing}
\renewcommand{\d}{\mathrm{d}}

\newcommand{\ra}{"$\Rightarrow$"}
\newcommand{\la}{"$\Leftarrow$"}
\newcommand{\bs}{\backslash}
\newcommand{\ind}{\large{1}}
\newcommand{\indic}{{\large 1}}

\DeclareMathAlphabet\mathbfcal{OMS}{cmsy}{b}{n}

% -------------------------------------------------
% Theorem styles
% -------------------------------------------------
\newtheoremstyle{mystyle}
  {}{}{}{}{
  \sffamily\bfseries}{.}{ }{}

\newtheoremstyle{cstyle}
  {}{}{}{}{
  \sffamily\bfseries}{.}{ }{\thmnote{#3}}

\makeatletter


\theoremstyle{mystyle}
\newtheorem{de}{Definition}[section]
\newtheorem{prop}[de]{Proposition}
\newtheorem{thm}[de]{Theorem}
\newtheorem{lem}[de]{Lemma}
\newtheorem{cor}[de]{Corollary}
\newtheorem*{remark}{Remark}
\newtheorem*{eg}{Example}

\theoremstyle{cstyle}
\newtheorem*{cthm}{}

\newtheoremstyle{warn}
  {}{}{}{}{\normalfont}{}{ }{}
\theoremstyle{warn}
\newtheorem*{warning}{\warningsign{0.2}\relax}

% -------------------------------------------------
% Warning symbol
% -------------------------------------------------
\newcommand{\warningsign}[1]{%
  \tikz[scale=#1,every node/.style={transform shape}]{%
    \draw[-,line width={#1*0.8mm},red,fill=yellow,
      rounded corners={#1*2.5mm}]
      (0,0)--(1,{-sqrt(3)})--(-1,{-sqrt(3)})--cycle;
    \node at (0,-1) {\fontsize{48}{60}\selectfont\bfseries!};%
  }%
}

% -------------------------------------------------
% Theorem boxes
% -------------------------------------------------
\tcolorboxenvironment{de}{
  boxrule=0pt,
  boxsep=0pt,
  colback=red!10,
  left=8pt,
  right=8pt,
  enhanced jigsaw,
  borderline west={2pt}{0pt}{red},
  sharp corners,
  before skip=10pt,
  after skip=10pt,
  breakable
}

\tcolorboxenvironment{prop}{
  boxrule=0pt,
  boxsep=0pt,
  colback=Orange!10,
  left=8pt,
  right=8pt,
  enhanced jigsaw,
  borderline west={2pt}{0pt}{Orange},
  sharp corners,
  before skip=10pt,
  after skip=10pt,
  breakable
}

\tcolorboxenvironment{thm}{
  boxrule=0pt,
  boxsep=0pt,
  colback=blue!10,
  left=8pt,
  right=8pt,
  enhanced jigsaw,
  borderline west={2pt}{0pt}{blue},
  sharp corners,
  before skip=10pt,
  after skip=10pt,
  breakable
}

\tcolorboxenvironment{lem}{
  boxrule=0pt,
  boxsep=0pt,
  colback=Cyan!10,
  left=8pt,
  right=8pt,
  enhanced jigsaw,
  borderline west={2pt}{0pt}{Cyan},
  sharp corners,
  before skip=10pt,
  after skip=10pt,
  breakable
}

\tcolorboxenvironment{cor}{
  boxrule=0pt,
  boxsep=0pt,
  colback=violet!10,
  left=8pt,
  right=8pt,
  enhanced jigsaw,
  borderline west={2pt}{0pt}{violet},
  sharp corners,
  before skip=10pt,
  after skip=10pt,
  breakable
}

\tcolorboxenvironment{proof}{
  boxrule=0pt,
  boxsep=0pt,
  blanker,
  borderline west={2pt}{0pt}{CadetBlue!80!white},
  left=8pt,
  right=8pt,
  sharp corners,
  before skip=10pt,
  after skip=10pt,
  breakable
}

\tcolorboxenvironment{remark}{
  boxrule=0pt,
  boxsep=0pt,
  blanker,
  borderline west={2pt}{0pt}{Green},
  left=8pt,
  right=8pt,
  before skip=10pt,
  after skip=10pt,
  breakable
}

\tcolorboxenvironment{eg}{
  boxrule=0pt,
  boxsep=0pt,
  blanker,
  borderline west={2pt}{0pt}{Black},
  left=8pt,
  right=8pt,
  sharp corners,
  before skip=10pt,
  after skip=10pt,
  breakable
}

\tcolorboxenvironment{cthm}{
  boxrule=0pt,
  boxsep=0pt,
  colback=gray!10,
  left=8pt,
  right=8pt,
  enhanced jigsaw,
  borderline west={2pt}{0pt}{gray},
  sharp corners,
  before skip=10pt,
  after skip=10pt,
  breakable
}

\newcommand{\for}{\operatorname{For}}
\newcommand{\fut}{\operatorname{Fut}}
\begin{document}
	\maketitle
	\allowdisplaybreaks
	\pagenumbering{roman}
	\setcounter{page}{1}
	\begin{center}
		\Large\textbf{Preface}
    \end{center}~\
    This is the note of the course \textit{Advanced financial models} taught in 2023 Spring by Prof. Gongqiu Zhang at the Chinese University of Hong Kong, Shenzhen, which covers a large fraction of the content in the book \textit{Stochastic Calculus in Finance II} by Steven Shreve. Most of the content is based on the presentation of the instructor and written down in my understanding, so there are unavoidable mistakes and typos in the note. The note is later revised and updated using ChatGPT. The note presumes basic knowledge in measure-based probability theory and discrete-time martingales, so the related contents in the lectures and the book are skipped or just stated without proof. Also some technical proofs in the book are also removed. Topics of the note include: basic probability and stochastic calculus, risk-neutral pricing theory, pricing of exotic options under Black-Scholes model, brief introductions to term-structure models and jump models.
	
	
	\newpage
	\tableofcontents
	\clearpage
    \pagenumbering{arabic}
\section{Review of probability and stochastic calculus}

\subsection{Brownian motions}

\begin{de}
A stochastic process $\{W_t\}_{t\geq 0}$ is called a standard Brownian motion if:
\begin{enumerate}
    \item $W_0=0$ almost surely;
    \item the map $t\mapsto W_t$ is continuous almost surely;
    \item $W$ has independent increments: for every $0\leq t_0<t_1<\cdots<t_n$, the random variables
    \[
    W_{t_1}-W_{t_0},\ldots,W_{t_n}-W_{t_{n-1}}
    \]
    are independent; and
    \item for every $0\leq s<t$,
    \[
    W_t-W_s\sim N(0,t-s).
    \]
\end{enumerate}
\end{de}

Brownian motion can be obtained as a scaling limit of a symmetric random walk. Let
\[
M_k=\sum_{i=1}^k X_i,
\qquad
\mathbb{P}(X_i=-1)=\mathbb{P}(X_i=1)=\frac12,
\]
and define
\[
W_t^{(n)}=\frac{1}{\sqrt n}M_{\lfloor nt\rfloor}.
\]
Donsker's invariance principle states that, after linear interpolation, $W^{(n)}$ converges in distribution to Brownian motion in the space of continuous paths on every compact time interval. This is convergence in law, not a pathwise limit on the original probability space.

\begin{thm}
Almost every Brownian path has the following properties.
\begin{enumerate}
    \item For every $T>0$ and every $\gamma<\frac12$, there is a finite random constant $C_{\gamma,T}$ such that
    \[
    |W_t-W_s|\leq C_{\gamma,T}|t-s|^\gamma,
    \qquad 0\leq s,t\leq T.
    \]
    Thus Brownian paths are locally H\"older continuous of every order strictly smaller than $\frac12$.
    \item The zero set
    \[
    \mathcal{Z}(\omega)=\{t\geq 0:W_t(\omega)=0\}
    \]
    is closed, unbounded, and has no isolated points. In particular, zeros and both positive and negative values occur arbitrarily close to time $0$.
    \item The strong law and the law of the iterated logarithm give
    \[
    \lim_{t\to\infty}\frac{W_t}{t}=0,
    \]
    and
    \[
    \limsup_{t\downarrow 0}
    \frac{W_t}{\sqrt{2t\log\log(1/t)}}=1,
    \qquad
    \liminf_{t\downarrow 0}
    \frac{W_t}{\sqrt{2t\log\log(1/t)}}=-1.
    \]
\end{enumerate}
\end{thm}

\begin{de}
Let $\Pi=\{0=t_0<t_1<\cdots<t_n=t\}$ be a partition of $[0,t]$, and let $X$ be a stochastic process. The $p$-variation sum of $X$ over $\Pi$ is
\[
V_t^{(p)}(\Pi,X)=\sum_{k=1}^n |X_{t_k}-X_{t_{k-1}}|^p.
\]
When these sums converge along a specified sequence of partitions whose mesh
\[
\|\Pi\|=\max_{1\leq k\leq n}(t_k-t_{k-1})
\]
tends to zero, the limit is called the $p$-variation along that sequence. For $p=2$, the limit is the quadratic variation and is denoted by $[X]_t=[X,X]_t$. For $p=1$, one obtains the total-variation sums; the total variation is the supremum of these sums over all partitions.
\end{de}

\begin{thm}
Let $W$ be a Brownian motion. Along any nested sequence of partitions of $[0,t]$ whose mesh tends to zero,
\[
[W]_t=t
\]
almost surely. Brownian paths have infinite total variation on every nontrivial interval and are nowhere differentiable almost surely.
\end{thm}

\begin{de}
Let $X$ and $Y$ be stochastic processes, and let $\Pi=\{0=t_0<t_1<\cdots<t_n=t\}$. If the limit exists, their quadratic covariation is
\[
[X,Y]_t
=
\lim_{\|\Pi\|\to 0}
\sum_{i=1}^n
(X_{t_i}-X_{t_{i-1}})(Y_{t_i}-Y_{t_{i-1}}).
\]
The mode of convergence and the sequence of partitions are understood from context.
\end{de}

\begin{thm}
Let $W$ be a Brownian motion. Then
\[
[W,t]_t=0,
\qquad
[t,t]_t=0.
\]
If $W^{(1)}$ and $W^{(2)}$ are independent Brownian motions, then
\[
[W^{(1)},W^{(2)}]_t=0.
\]
In It\^o differential notation,
\[
(dW_t)^2=dt,
\qquad
dW_t\,dt=0,
\qquad
(dt)^2=0,
\qquad
dW_t^{(1)}\,dW_t^{(2)}=0.
\]
More generally, for Brownian motions with instantaneous correlation $\rho_t$,
\[
dW_t^{(1)}\,dW_t^{(2)}=\rho_t\,dt.
\]
\end{thm}

\begin{de}
Let $W$ be a Brownian motion on $(\Omega,\mathcal{F},\mathbb{P})$. A filtration $\{\mathcal{F}_t\}_{t\geq 0}$ is a Brownian filtration for $W$ if:
\begin{enumerate}
    \item $W_t$ is $\mathcal{F}_t$-measurable for every $t\geq 0$; and
    \item for every $0\leq s<t$, the increment $W_t-W_s$ is independent of $\mathcal{F}_s$.
\end{enumerate}
Usually the filtration is also assumed to satisfy the usual conditions of completeness and right-continuity.
\end{de}

\begin{remark}
The natural filtration of $W$ is
\[
\mathcal{F}_t^W=\sigma(W_s:0\leq s\leq t),
\]
usually completed and made right-continuous. A Brownian motion may also be adapted to larger filtrations, provided its future increments remain independent of the past.
\end{remark}

\begin{thm}
Let $W$ be a standard Brownian motion with Brownian filtration $\{\mathcal{F}_t\}$. Then:
\begin{enumerate}
    \item for every $a>0$, the process $a^{-1/2}W_{at}$ is a standard Brownian motion;
    \item
    \[
    \operatorname{Cov}(W_s,W_t)=s\wedge t;
    \]
    \item $W$ is a martingale:
    \[
    \mathbb{E}[W_t\mid\mathcal{F}_s]=W_s,
    \qquad 0\leq s\leq t;
    \]
    \item $W$ is Markov: for every bounded Borel function $h$,
    \[
    \mathbb{E}[h(W_t)\mid\mathcal{F}_s]
    =
    (P_{t-s}h)(W_s),
    \]
    where
    \[
    (P_u h)(x)=\mathbb{E}[h(x+\sqrt{u}\,Z)]
    \quad\text{and}\quad Z\sim N(0,1).
    \]
\end{enumerate}
\end{thm}

\subsection{More on Brownian motions}

\begin{thm}
Let $0<t_1<\cdots<t_m$, and let $W$ be a standard Brownian motion. The moment-generating function of
\[
(W_{t_1},\ldots,W_{t_m})
\]
is
\[
\begin{aligned}
\varphi(u_1,\ldots,u_m)
&=
\mathbb{E}\exp\left(\sum_{j=1}^m u_jW_{t_j}\right)\\
&=
\exp\left\{
\frac12\sum_{k=1}^m
\left(\sum_{j=k}^m u_j\right)^2(t_k-t_{k-1})
\right\},
\end{aligned}
\]
where $t_0=0$.
\end{thm}

\begin{proof}
Write
\[
\sum_{j=1}^m u_jW_{t_j}
=
\sum_{k=1}^m
\left(\sum_{j=k}^m u_j\right)(W_{t_k}-W_{t_{k-1}}).
\]
The increments are independent and Gaussian, so their moment-generating functions multiply. Since
\[
W_{t_k}-W_{t_{k-1}}\sim N(0,t_k-t_{k-1}),
\]
the displayed formula follows.
\end{proof}

\begin{thm}[Exponential martingale]
Let $W$ be a Brownian motion and let $\sigma\in\mathbb{R}$ be constant. Then
\[
Z_t=\exp\left(\sigma W_t-\frac12\sigma^2t\right),
\qquad t\geq 0,
\]
is a martingale.
\end{thm}

\begin{proof}
For $0\leq s\leq t$, independence of $W_t-W_s$ and $\mathcal{F}_s$ gives
\[
\begin{aligned}
\mathbb{E}[Z_t\mid\mathcal{F}_s]
&=
\exp\left(\sigma W_s-\frac12\sigma^2t\right)
\mathbb{E}\left[e^{\sigma(W_t-W_s)}\right]\\
&=
\exp\left(\sigma W_s-\frac12\sigma^2t\right)
\exp\left(\frac12\sigma^2(t-s)\right)\\
&=Z_s.
\end{aligned}
\]
\end{proof}

\begin{de}
For $m\in\mathbb{R}$, define the first passage time to level $m$ by
\[
\tau_m=\inf\{t\geq 0:W_t=m\},
\]
with the convention $\inf\varnothing=\infty$.
\end{de}

\begin{thm}
For every $m\in\mathbb{R}$, $\tau_m<\infty$ almost surely. Moreover, for every $\alpha>0$,
\[
\mathbb{E}\left[e^{-\alpha\tau_m}\right]
=e^{-|m|\sqrt{2\alpha}}.
\]
For $m=0$, this is interpreted using $\tau_0=0$.
\end{thm}

\begin{proof}
It is enough by symmetry to consider $m>0$. Let $\sigma>0$ and stop the exponential martingale at $t\wedge\tau_m$. Since $t\wedge\tau_m$ is bounded,
\[
1
=
\mathbb{E}\left[
\exp\left(
\sigma W_{t\wedge\tau_m}
-\frac12\sigma^2(t\wedge\tau_m)
\right)
\right].
\]
Before $\tau_m$, continuity implies $W_{t\wedge\tau_m}\leq m$. Dominated convergence therefore gives, as $t\to\infty$,
\[
1
=
e^{\sigma m}
\mathbb{E}\left[
\mathbf{1}_{\{\tau_m<\infty\}}
e^{-\frac12\sigma^2\tau_m}
\right].
\]
Letting $\sigma\downarrow 0$ yields $\mathbb{P}(\tau_m<\infty)=1$. Hence
\[
\mathbb{E}\left[e^{-\frac12\sigma^2\tau_m}\right]
=e^{-\sigma m}.
\]
Putting $\alpha=\sigma^2/2$ gives the stated formula for $m>0$, and symmetry gives the result for $m<0$.
\end{proof}

\begin{remark}
Differentiating the Laplace transform for $m\neq 0$ gives
\[
\mathbb{E}\left[\tau_m e^{-\alpha\tau_m}\right]
=
\frac{|m|}{\sqrt{2\alpha}}
 e^{-|m|\sqrt{2\alpha}},
\qquad \alpha>0.
\]
Letting $\alpha\downarrow 0$ shows that $\mathbb{E}[\tau_m]=\infty$ for $m\neq 0$.
\end{remark}

\begin{thm}[Reflection principle]
Let $m>0$ and $w\leq m$. Then
\[
\mathbb{P}(\tau_m\leq t,\,W_t\leq w)
=
\mathbb{P}(W_t\geq 2m-w).
\]
\end{thm}

\begin{proof}
Reflect each path after its first hitting time of $m$. This gives a measure-preserving bijection between paths satisfying $\tau_m\leq t$ and $W_t\leq w$ and paths satisfying $W_t\geq 2m-w$.
\end{proof}

\begin{thm}
For $m\neq 0$ and $t>0$,
\[
\mathbb{P}(\tau_m\leq t)
=
2\left[1-\Phi\left(\frac{|m|}{\sqrt t}\right)\right]
=
\frac{2}{\sqrt{2\pi}}
\int_{|m|/\sqrt t}^{\infty}e^{-y^2/2}\,dy,
\]
and $\tau_m$ has density
\[
f_{\tau_m}(t)
=
\frac{|m|}{\sqrt{2\pi t^3}}
\exp\left(-\frac{m^2}{2t}\right),
\qquad t>0.
\]
\end{thm}

\begin{proof}
For $m>0$, the reflection principle gives
\[
\mathbb{P}(\tau_m\leq t)
=2\mathbb{P}(W_t\geq m).
\]
The formula follows from $W_t\sim N(0,t)$. The case $m<0$ follows by symmetry, and differentiation gives the density.
\end{proof}

\begin{thm}
Let
\[
M_t=\max_{0\leq s\leq t}W_s.
\]
For $m>0$ and $w\leq m$,
\[
\mathbb{P}(M_t\geq m,\,W_t\leq w)
=
\mathbb{P}(W_t\geq 2m-w).
\]
The joint density of $(M_t,W_t)$ is
\[
f_{M_t,W_t}(m,w)
=
\frac{2(2m-w)}{t\sqrt{2\pi t}}
\exp\left(-\frac{(2m-w)^2}{2t}\right),
\qquad m\geq 0,\quad w\leq m,
\]
and is zero elsewhere.
\end{thm}

\begin{proof}
The reflection identity gives
\[
\int_m^\infty\int_{-\infty}^w
f_{M_t,W_t}(x,y)\,dy\,dx
=
\frac{1}{\sqrt{2\pi t}}
\int_{2m-w}^{\infty}e^{-z^2/(2t)}\,dz.
\]
Differentiate first with respect to $m$ and then with respect to $w$.
\end{proof}

\subsection{It\^o integrals}

Let $\mathcal{H}^2([0,T])$ denote the space of progressively measurable processes $X$ such that
\[
\|X\|_{\mathcal{H}^2}
=
\left(
\mathbb{E}\int_0^T |X_s|^2\,ds
\right)^{1/2}
<\infty.
\]

\begin{de}
An adapted simple process on $[0,T]$ has the form
\[
\phi_s
=
\sum_{i=0}^{n-1}
A_i\mathbf{1}_{(t_i,t_{i+1}]}(s),
\]
where
\[
0=t_0<t_1<\cdots<t_n=T,
\qquad
A_i\in L^2(\Omega,\mathcal{F}_{t_i},\mathbb{P}).
\]
\end{de}

\begin{de}
For an adapted simple process $\phi$, define
\[
\int_0^t \phi_s\,dW_s
=
\sum_{i=0}^{n-1}
A_i\bigl(W_{t\wedge t_{i+1}}-W_{t\wedge t_i}\bigr),
\qquad 0\leq t\leq T.
\]
Equivalently, if $t\in(t_k,t_{k+1}]$, then
\[
\int_0^t \phi_s\,dW_s
=
\sum_{i=0}^{k-1}A_i(W_{t_{i+1}}-W_{t_i})
+A_k(W_t-W_{t_k}).
\]
\end{de}

\begin{thm}[Density of simple processes]
For every $X\in\mathcal{H}^2([0,T])$, there is a sequence of adapted simple processes $\phi^{(n)}$ such that
\[
\lim_{n\to\infty}
\mathbb{E}\int_0^T
|\phi_s^{(n)}-X_s|^2\,ds
=0.
\]
\end{thm}

\begin{de}
For $X\in\mathcal{H}^2([0,T])$, choose adapted simple processes $\phi^{(n)}$ as above and define
\[
\int_0^t X_s\,dW_s
=
L^2(\Omega)\text{-}\lim_{n\to\infty}
\int_0^t\phi_s^{(n)}\,dW_s.
\]
The It\^o isometry shows that this limit is independent of the approximating sequence.
\end{de}

\begin{thm}
Let $X,Y\in\mathcal{H}^2([0,T])$ and $a,b\in\mathbb{R}$.
\begin{enumerate}
    \item The It\^o integral is linear:
    \[
    \int_0^t(aX_s+bY_s)\,dW_s
    =
    a\int_0^tX_s\,dW_s
    +b\int_0^tY_s\,dW_s.
    \]
    \item The process
    \[
    I_t=\int_0^tX_s\,dW_s
    \]
    is a square-integrable martingale with $\mathbb{E}[I_t]=0$.
    \item The It\^o isometry holds:
    \[
    \mathbb{E}\left[
    \left(\int_0^tX_s\,dW_s\right)^2
    \right]
    =
    \mathbb{E}\int_0^tX_s^2\,ds.
    \]
    \item The quadratic variation is
    \[
    [I]_t=\int_0^tX_s^2\,ds,
    \]
    or, in differential notation,
    \[
    (dI_t)^2=X_t^2\,dt.
    \]
\end{enumerate}
\end{thm}

\begin{de}
An It\^o process is a process of the form
\[
X_t
=
X_0+
\int_0^t b_s\,ds+
\int_0^t\sigma_s\,dW_s,
\]
where the integrals are well defined. In differential notation,
\[
dX_t=b_t\,dt+\sigma_t\,dW_t.
\]
Its quadratic variation is
\[
[X]_t=\int_0^t\sigma_s^2\,ds.
\]
\end{de}

\begin{de}
If $X$ is the It\^o process above and $Y$ is progressively measurable with
\[
\int_0^t|Y_sb_s|\,ds<\infty
\quad\text{and}\quad
\mathbb{E}\int_0^tY_s^2\sigma_s^2\,ds<\infty,
\]
define
\[
\int_0^tY_s\,dX_s
=
\int_0^tY_sb_s\,ds+
\int_0^tY_s\sigma_s\,dW_s.
\]
\end{de}

\begin{thm}[It\^o's formula]
Let
\[
dX_t=b_t\,dt+\sigma_t\,dW_t,
\]
and let $f\in C^{1,2}([0,T]\times\mathbb{R})$. Then
\[
\begin{aligned}
f(t,X_t)
={}&f(0,X_0)
+\int_0^t f_t(s,X_s)\,ds
+\int_0^t f_x(s,X_s)\,dX_s\\
&+\frac12\int_0^t f_{xx}(s,X_s)\,d[X]_s.
\end{aligned}
\]
Equivalently,
\[
df(t,X_t)
=
f_t(t,X_t)\,dt
+f_x(t,X_t)\,dX_t
+\frac12f_{xx}(t,X_t)(dX_t)^2.
\]
\end{thm}

\begin{thm}[Integration by parts]
Let $f:[0,T]\to\mathbb{R}$ be deterministic, continuous, and of bounded variation. Then
\[
f(t)W_t-f(0)W_0
=
\int_0^t f(s)\,dW_s+
\int_0^t W_s\,df(s).
\]
Since $W_0=0$,
\[
\int_0^t f(s)\,dW_s
=
f(t)W_t-
\int_0^t W_s\,df(s).
\]
\end{thm}

\begin{eg}
Let
\[
Y_t
=
\exp\left(
-\frac12\int_0^t\theta_s^2\,ds
-\int_0^t\theta_s\,dW_s
\right).
\]
Set
\[
X_t
=-\frac12\int_0^t\theta_s^2\,ds
-\int_0^t\theta_s\,dW_s.
\]
Then $Y_t=e^{X_t}$ and It\^o's formula gives
\[
\begin{aligned}
dY_t
&=e^{X_t}\,dX_t+\frac12e^{X_t}(dX_t)^2\\
&=e^{X_t}\left(-\frac12\theta_t^2\,dt-\theta_t\,dW_t\right)
+\frac12e^{X_t}\theta_t^2\,dt\\
&=-Y_t\theta_t\,dW_t.
\end{aligned}
\]
Hence
\[
Y_t=1-\int_0^tY_s\theta_s\,dW_s.
\]
Under an additional condition such as Novikov's condition, $Y$ is a true martingale; without such a condition, it is always a nonnegative local martingale.
\end{eg}

\begin{thm}[L\'evy's characterization]
Let $M$ be a continuous local martingale such that
\[
M_0=0
\quad\text{and}\quad
[M]_t=t.
\]
Then $M$ is a standard Brownian motion with respect to the underlying filtration.
\end{thm}

\begin{de}
A $d$-dimensional process
\[
W_t=(W_t^{(1)},\ldots,W_t^{(d)})^\top
\]
is a standard $d$-dimensional Brownian motion if it has continuous paths, $W_0=0$, independent increments, and
\[
W_t-W_s\sim N(0,(t-s)I_d)
\]
for $0\leq s<t$.
\end{de}

\begin{remark}
A standard $d$-dimensional Brownian motion can be constructed by collecting $d$ independent one-dimensional Brownian motions into a vector.
\end{remark}

\begin{thm}[Multidimensional L\'evy characterization]
Let $M=(M^{(1)},\ldots,M^{(d)})^\top$ be a continuous $d$-dimensional local martingale with $M_0=0$. If
\[
[M^{(i)},M^{(j)}]_t=t\delta_{ij}
\]
for every $i,j$, then $M$ is a standard $d$-dimensional Brownian motion.
\end{thm}

\begin{de}
An $n$-dimensional It\^o process driven by a $d$-dimensional Brownian motion is a process
\[
X_t
=
X_0+
\int_0^t b_s\,ds+
\int_0^t\Sigma_s\,dW_s,
\]
where $b_s\in\mathbb{R}^n$ and $\Sigma_s\in\mathbb{R}^{n\times d}$ are progressively measurable and satisfy the required integrability conditions.
\end{de}

\begin{thm}[Multidimensional It\^o formula]
Let $X$ be the $n$-dimensional It\^o process above, and let $f\in C^{1,2}([0,T]\times\mathbb{R}^n)$. Then
\[
\begin{aligned}
f(t,X_t)
={}&f(0,X_0)
+\int_0^t f_t(s,X_s)\,ds
+\int_0^t\nabla f(s,X_s)^\top\,dX_s\\
&+\frac12\sum_{i,j=1}^n
\int_0^t
f_{x_ix_j}(s,X_s)\,d[X^{(i)},X^{(j)}]_s.
\end{aligned}
\]
In differential notation,
\[
df(t,X_t)
=
f_t(t,X_t)\,dt
+\nabla f(t,X_t)^\top dX_t
+\frac12\sum_{i,j=1}^n
f_{x_ix_j}(t,X_t)\,d[X^{(i)},X^{(j)}]_t.
\]
If $dX_t=b_t\,dt+\Sigma_t\,dW_t$, then
\[
d[X^{(i)},X^{(j)}]_t
=(\Sigma_t\Sigma_t^\top)_{ij}\,dt.
\]
\end{thm}

\begin{thm}[Product rule]
For two continuous semimartingales $X$ and $Y$,
\[
d(X_tY_t)
=X_t\,dY_t+Y_t\,dX_t+d[X,Y]_t.
\]
For It\^o processes, the last term is often written $dX_t\,dY_t$.
\end{thm}

\subsection{Brownian bridges}

\begin{de}
A stochastic process $X=\{X_t\}_{t\geq 0}$ is Gaussian if, for every finite collection of times $t_1,\ldots,t_n$, the random vector
\[
(X_{t_1},\ldots,X_{t_n})
\]
is jointly Gaussian.
\end{de}

\begin{eg}
Let $\Delta$ be a deterministic function satisfying
\[
\int_0^T\Delta(s)^2\,ds<\infty
\]
for every $T>0$. Then
\[
I_t=\int_0^t\Delta(s)\,dW_s
\]
is a Gaussian process.
\end{eg}

\begin{de}
Fix $T>0$. The Brownian bridge from $0$ to $0$ on $[0,T]$ is
\[
X_t=W_t-\frac{t}{T}W_T,
\qquad 0\leq t\leq T.
\]
\end{de}

\begin{thm}
The Brownian bridge $X$ is a continuous Gaussian process with
\[
\mathbb{E}[X_t]=0
\]
and covariance
\[
\operatorname{Cov}(X_s,X_t)
=s\wedge t-\frac{st}{T}.
\]
Moreover, $X_0=X_T=0$.
\end{thm}

\begin{proof}
Every finite vector of bridge values is a linear transformation of a finite Gaussian vector of Brownian values, so it is jointly Gaussian. Also,
\[
\begin{aligned}
\mathbb{E}[X_sX_t]
&=
\mathbb{E}\left[
\left(W_s-\frac{s}{T}W_T\right)
\left(W_t-\frac{t}{T}W_T\right)
\right]\\
&=s\wedge t-\frac{st}{T}.
\end{aligned}
\]
Continuity follows from the continuity of $W$.
\end{proof}

\begin{remark}
The process $X_t=W_t-(t/T)W_T$ is not adapted to the natural filtration of $W$ because it uses the future value $W_T$.
\end{remark}

\begin{de}
For $a,b\in\mathbb{R}$, the Brownian bridge from $a$ to $b$ on $[0,T]$ is
\[
X_t^{a\to b}
=
a+\frac{b-a}{T}t+X_t.
\]
Its mean and covariance are
\[
\mathbb{E}[X_t^{a\to b}]
=a+\frac{b-a}{T}t,
\qquad
\operatorname{Cov}(X_s^{a\to b},X_t^{a\to b})
=s\wedge t-\frac{st}{T}.
\]
\end{de}

An adapted version with the same law is
\[
Y_t
=
(T-t)\int_0^t\frac{1}{T-u}\,dW_u,
\qquad 0\leq t<T.
\]
For $0\leq s\leq t<T$,
\[
\begin{aligned}
\operatorname{Cov}(Y_s,Y_t)
&=(T-s)(T-t)
\int_0^s\frac{1}{(T-u)^2}\,du\\
&=(T-s)(T-t)
\left(\frac{1}{T-s}-\frac{1}{T}\right)\\
&=s-\frac{st}{T}.
\end{aligned}
\]
Thus $Y$ has the same finite-dimensional distributions as the Brownian bridge.

\begin{thm}
Define
\[
Y_t=
\begin{cases}
(T-t)\displaystyle\int_0^t\frac{1}{T-u}\,dW_u,
&0\leq t<T,\\[1.2ex]
0,&t=T.
\end{cases}
\]
Then $Y$ has a continuous version on $[0,T]$ and has the same law as the Brownian bridge from $0$ to $0$.
\end{thm}

\begin{proof}
The mean and covariance calculations above identify the finite-dimensional distributions. To verify continuity at $T$, note that the quadratic variation of
\[
I_t=\int_0^t\frac{1}{T-u}\,dW_u
\]
is
\[
q(t)=\frac{1}{T-t}-\frac1T.
\]
By the Dambis--Dubins--Schwarz time-change representation, $I_t=\widehat W_{q(t)}$ for a Brownian motion $\widehat W$. Since $\widehat W_s/s\to0$ almost surely as $s\to\infty$ and $(T-t)q(t)\to1$, it follows that $(T-t)I_t\to0$ almost surely as $t\uparrow T$.
\end{proof}

\begin{remark}
For $t<T$, the adapted bridge satisfies
\[
dY_t=-\frac{Y_t}{T-t}\,dt+dW_t.
\]
The singular mean-reversion term forces the process toward zero as $t\uparrow T$.
\end{remark}

\subsection{Change of measure by random variables}

\begin{thm}
Let $Z$ be a nonnegative $\mathcal{F}$-measurable random variable on $(\Omega,\mathcal{F},\mathbb{P})$ with $\mathbb{E}_{\mathbb{P}}[Z]=1$. Then
\[
\widetilde{\mathbb{P}}(A)
=
\mathbb{E}_{\mathbb{P}}[Z\mathbf{1}_A]
=
\int_A Z\,d\mathbb{P},
\qquad A\in\mathcal{F},
\]
defines a probability measure.
\end{thm}

\begin{de}
If $\widetilde{\mathbb{P}}\ll\mathbb{P}$, the Radon--Nikodym derivative
\[
Z=\frac{d\widetilde{\mathbb{P}}}{d\mathbb{P}}
\]
is the nonnegative random variable satisfying
\[
\widetilde{\mathbb{P}}(A)=\int_A Z\,d\mathbb{P}
\]
for every $A\in\mathcal{F}$.
\end{de}

\begin{thm}
If $X$ is integrable under $\widetilde{\mathbb{P}}$, then
\[
\widetilde{\mathbb{E}}[X]
=
\mathbb{E}[ZX].
\]
If $Z>0$ almost surely, so that $\widetilde{\mathbb{P}}$ and $\mathbb{P}$ are equivalent, then for every $\mathbb{P}$-integrable $X$,
\[
\mathbb{E}[X]
=
\widetilde{\mathbb{E}}\left[\frac{X}{Z}\right].
\]
\end{thm}

\begin{de}
Two probability measures $\mathbb{P}$ and $\widetilde{\mathbb{P}}$ on $(\Omega,\mathcal{F})$ are equivalent, written
\[
\widetilde{\mathbb{P}}\sim\mathbb{P},
\]
if they have the same null sets.
\end{de}

\begin{eg}
Let $X\sim N(0,1)$ under $\mathbb{P}$ and set
\[
Z=e^{-\theta X-\frac12\theta^2}.
\]
Then $\mathbb{E}[Z]=1$, and the probability measure defined by $Z$ satisfies
\[
\begin{aligned}
\widetilde{\mathbb{E}}[e^{iuX}]
&=
\mathbb{E}[Ze^{iuX}]\\
&=
\exp\left(-\frac12u^2-iu\theta\right).
\end{aligned}
\]
Hence $X\sim N(-\theta,1)$ under $\widetilde{\mathbb{P}}$.
\end{eg}

\subsection{Change of measure by stochastic processes}

Fix a terminal time $T$. Suppose $\widetilde{\mathbb{P}}\ll\mathbb{P}$ on $\mathcal{F}_T$, and define
\[
Z_T=\frac{d\widetilde{\mathbb{P}}}{d\mathbb{P}}
\quad\text{and}\quad
Z_t=\mathbb{E}_{\mathbb{P}}[Z_T\mid\mathcal{F}_t],
\qquad 0\leq t\leq T.
\]
Then $Z$ is a nonnegative $\mathbb{P}$-martingale with $\mathbb{E}[Z_t]=1$, and
\[
Z_t
=
\left.
\frac{d\widetilde{\mathbb{P}}}{d\mathbb{P}}
\right|_{\mathcal{F}_t}.
\]

\begin{thm}[Bayes formula]
Assume $\widetilde{\mathbb{P}}\sim\mathbb{P}$, so $Z_t>0$ almost surely. If $0\leq s\leq t\leq T$ and $Y$ is $\mathcal{F}_t$-measurable and integrable under $\widetilde{\mathbb{P}}$, then
\[
\widetilde{\mathbb{E}}[Y\mid\mathcal{F}_s]
=
\frac{1}{Z_s}
\mathbb{E}[Z_tY\mid\mathcal{F}_s].
\]
\end{thm}

\begin{proof}
The right-hand side is $\mathcal{F}_s$-measurable. For every $A\in\mathcal{F}_s$,
\[
\begin{aligned}
\widetilde{\mathbb{E}}\left[
\mathbf{1}_A
\frac{1}{Z_s}\mathbb{E}[Z_tY\mid\mathcal{F}_s]
\right]
&=
\mathbb{E}\left[
Z_s\mathbf{1}_A
\frac{1}{Z_s}\mathbb{E}[Z_tY\mid\mathcal{F}_s]
\right]\\
&=
\mathbb{E}[\mathbf{1}_AZ_tY]\\
&=
\widetilde{\mathbb{E}}[\mathbf{1}_AY].
\end{aligned}
\]
\end{proof}

\begin{thm}[Girsanov's theorem]
Let $W$ be a Brownian motion under $\mathbb{P}$, and let $\theta$ be progressively measurable. Assume Novikov's condition
\[
\mathbb{E}\exp\left(
\frac12\int_0^T\theta_s^2\,ds
\right)<\infty.
\]
Define
\[
Z_t
=
\exp\left(
-\int_0^t\theta_s\,dW_s
-\frac12\int_0^t\theta_s^2\,ds
\right).
\]
Then $Z$ is a strictly positive martingale. Define $\widetilde{\mathbb{P}}$ on $\mathcal{F}_T$ by
\[
\frac{d\widetilde{\mathbb{P}}}{d\mathbb{P}}=Z_T.
\]
The process
\[
\widetilde W_t
=W_t+\int_0^t\theta_s\,ds
\]
is a Brownian motion under $\widetilde{\mathbb{P}}$.
\end{thm}

\begin{proof}
It\^o's formula gives
\[
dZ_t=-Z_t\theta_t\,dW_t.
\]
Novikov's condition upgrades the stochastic exponential from a local martingale to a true martingale. Moreover,
\[
\begin{aligned}
d(Z_t\widetilde W_t)
&=Z_t\,d\widetilde W_t+
\widetilde W_t\,dZ_t+d[Z,\widetilde W]_t\\
&=Z_t(dW_t+\theta_t\,dt)
-\widetilde W_tZ_t\theta_t\,dW_t
-Z_t\theta_t\,dt\\
&=Z_t(1-\widetilde W_t\theta_t)\,dW_t.
\end{aligned}
\]
After localization, $Z_t\widetilde W_t$ is a $\mathbb{P}$-martingale. Bayes' formula therefore shows that $\widetilde W$ is a $\widetilde{\mathbb{P}}$-local martingale. Its quadratic variation is
\[
[\widetilde W]_t=[W]_t=t,
\]
so L\'evy's characterization implies that $\widetilde W$ is a Brownian motion under $\widetilde{\mathbb{P}}$.
\end{proof}

\subsection{Stochastic differential equations}

\begin{de}
Given $t\leq u\leq T$ and $X_t=x$, a solution of
\[
dX_u=\beta(u,X_u)\,du+\gamma(u,X_u)\,dW_u
\]
satisfies the integral equation
\[
X_u
=
x+
\int_t^u\beta(s,X_s)\,ds+
\int_t^u\gamma(s,X_s)\,dW_s.
\]
\end{de}

\begin{remark}
The Euler--Maruyama approximation with step size $\Delta t$ is
\[
X_{t+\Delta t}
\approx
X_t+
\beta(t,X_t)\Delta t+
\gamma(t,X_t)(W_{t+\Delta t}-W_t).
\]
It can be used to simulate approximate sample paths of $X$.
\end{remark}

\begin{eg}[Geometric Brownian motion]
Suppose
\[
dS_t=\mu S_t\,dt+\sigma S_t\,dW_t,
\qquad S_0>0.
\]
It\^o's formula applied to $\log S_t$ gives
\[
d\log S_t
=
\left(\mu-\frac12\sigma^2\right)dt+
\sigma\,dW_t.
\]
Therefore,
\[
S_t
=
S_0\exp\left[
\left(\mu-\frac12\sigma^2\right)t+
\sigma W_t
\right].
\]
\end{eg}

\begin{eg}[Ornstein--Uhlenbeck process]
Suppose
\[
dR_t=\kappa(\theta-R_t)\,dt+\sigma\,dW_t.
\]
Multiplying by $e^{\kappa t}$ gives
\[
d(e^{\kappa t}R_t)
=
\kappa\theta e^{\kappa t}\,dt+
\sigma e^{\kappa t}\,dW_t.
\]
Hence
\[
R_t
=
\theta+(R_0-\theta)e^{-\kappa t}
+
\sigma\int_0^t e^{-\kappa(t-s)}\,dW_s.
\]
\end{eg}

\begin{eg}[Linear SDE]
Consider
\[
dX_t
=
(a(t)+b(t)X_t)\,dt+
(c(t)+h(t)X_t)\,dW_t.
\]
Define
\[
Y_t
=
\exp\left(
\int_0^t\left(b(s)-\frac12h(s)^2\right)ds
+
\int_0^t h(s)\,dW_s
\right).
\]
Then
\[
X_t
=
Y_t\left[
X_0+
\int_0^t\frac{a(s)-c(s)h(s)}{Y_s}\,ds
+
\int_0^t\frac{c(s)}{Y_s}\,dW_s
\right].
\]
\end{eg}

\begin{proof}
It\^o's formula gives
\[
dY_t=b(t)Y_t\,dt+h(t)Y_t\,dW_t.
\]
Set
\[
Z_t
=
X_0+
\int_0^t\frac{a(s)-c(s)h(s)}{Y_s}\,ds
+
\int_0^t\frac{c(s)}{Y_s}\,dW_s.
\]
Since $X_t=Y_tZ_t$, the product rule gives
\[
\begin{aligned}
dX_t
&=Y_t\,dZ_t+Z_t\,dY_t+d[Y,Z]_t\\
&=(a(t)+b(t)X_t)\,dt
+(c(t)+h(t)X_t)\,dW_t.
\end{aligned}
\]
\end{proof}

\begin{thm}[Markov property of SDE solutions]
Suppose $\beta(t,x)$ and $\gamma(t,x)$ are deterministic coefficients satisfying conditions that guarantee a unique strong solution, such as global Lipschitz and linear-growth conditions. Then the solution of
\[
dX_t=\beta(t,X_t)\,dt+\gamma(t,X_t)\,dW_t
\]
is a time-inhomogeneous Markov process.
\end{thm}

\begin{thm}[Feynman--Kac formula]
Let $X$ solve
\[
dX_s=\beta(s,X_s)\,ds+\gamma(s,X_s)\,dW_s,
\qquad X_t=x.
\]
Under standard regularity and integrability assumptions, define
\[
f(t,x)
=
\mathbb{E}_{t,x}\left[
\exp\left(-\int_t^T\kappa(u,X_u)\,du\right)
 h(X_T)
\right].
\]
If $f\in C^{1,2}$, then it solves the terminal-value problem
\[
\begin{cases}
 f_t+\beta f_x+\dfrac12\gamma^2 f_{xx}-\kappa f=0,
 &t<T,\\[0.8ex]
 f(T,x)=h(x).
\end{cases}
\]
\end{thm}

\begin{proof}
The process
\[
M_s
=
\exp\left(-\int_t^s\kappa(u,X_u)\,du\right)f(s,X_s),
\qquad t\leq s\leq T,
\]
is a martingale by the tower property. Applying the product rule and It\^o's formula gives
\[
\begin{aligned}
dM_s
={}&
\exp\left(-\int_t^s\kappa(u,X_u)\,du\right)
\left(
 f_t+\beta f_x+\frac12\gamma^2f_{xx}-\kappa f
\right)(s,X_s)\,ds\\
&+
\exp\left(-\int_t^s\kappa(u,X_u)\,du\right)
\gamma(s,X_s)f_x(s,X_s)\,dW_s.
\end{aligned}
\]
The drift must vanish, and the terminal condition follows by setting $t=T$.
\end{proof}

\begin{remark}
A common method for computing conditional expectations is to construct a martingale, apply It\^o's formula, and set the drift term equal to zero. This produces a partial differential equation.
\end{remark}

\begin{de}
A multidimensional SDE has the form
\[
dX_t=b(t,X_t)\,dt+\Sigma(t,X_t)\,dW_t,
\]
where $X_t\in\mathbb{R}^m$, $W_t\in\mathbb{R}^d$, $b:[0,T]\times\mathbb{R}^m\to\mathbb{R}^m$, and $\Sigma:[0,T]\times\mathbb{R}^m\to\mathbb{R}^{m\times d}$.
\end{de}

\begin{thm}[Multidimensional Feynman--Kac formula]
Let $X$ solve the multidimensional SDE above, and define
\[
f(t,x)
=
\mathbb{E}_{t,x}\left[
\exp\left(-\int_t^T\kappa(u,X_u)\,du\right)
 h(X_T)
\right].
\]
Under standard regularity assumptions, $f$ solves
\[
\begin{cases}
 f_t+b(t,x)^\top\nabla f
 +\dfrac12\operatorname{tr}\!\left(
 \Sigma(t,x)\Sigma(t,x)^\top D^2f
 \right)
 -\kappa(t,x)f=0,\\[1ex]
 f(T,x)=h(x).
\end{cases}
\]
Equivalently, the second-order term is
\[
\frac12\sum_{i,j=1}^m
(\Sigma\Sigma^\top)_{ij}(t,x)
\frac{\partial^2f}{\partial x_i\partial x_j}(t,x).
\]
\end{thm}

    \section{Risk-neutral pricing}

\subsection{Introduction}

The purpose of an arbitrage-pricing model is to describe the prices of traded assets and to assign values to contingent claims in a way that is consistent with the absence of arbitrage. The main financial instruments considered in these notes are the following.

\begin{itemize}
    \item \textbf{Stocks.} A share of stock represents partial ownership of a company. Shareholders may receive dividends, and the stock price reflects market expectations about the firm's future cash flows and risks.
    \item \textbf{Bonds.} A bond is a debt security issued by a corporation, government, or other entity. The issuer promises specified cash flows, such as coupons and principal repayment. Bond prices are sensitive to interest rates and credit risk.
    \item \textbf{Bank deposits and money-market accounts.} Money invested in a bank account earns interest. In an idealized frictionless model, the money-market account is the locally risk-free traded asset.
    \item \textbf{Commodities.} Commodities are physical goods such as oil, gold, and agricultural products. Their prices are affected by supply, demand, storage costs, convenience yields, and broader economic conditions.
    \item \textbf{Financial derivatives.} A derivative is a contract whose payoff depends on one or more underlying assets or market variables. Examples include forwards, futures, and options.
\end{itemize}

A central restriction on any pricing model is the absence of arbitrage.

\begin{de}
Let $X$ be the wealth process of an admissible self-financing trading strategy. The strategy is an \emph{arbitrage} on $[0,T]$ if
\[
X_0=0,
\qquad
\mathbb{P}(X_T\geq 0)=1,
\qquad
\mathbb{P}(X_T>0)>0.
\]
Admissibility is imposed to exclude doubling strategies; for example, one may require the discounted wealth process to be bounded from below.
\end{de}

An equivalent martingale measure is the main tool used to rule out arbitrage.

\begin{de}
Let $B$ be a strictly positive numeraire, usually the money-market account, and let $S^{(1)},\ldots,S^{(m)}$ be traded assets. A probability measure $\widetilde{\mathbb{P}}$ is an \emph{equivalent martingale measure} for the numeraire $B$ if
\begin{enumerate}
    \item $\widetilde{\mathbb{P}}\sim\mathbb{P}$; and
    \item every discounted traded-asset price $S_t^{(i)}/B_t$ is a $\widetilde{\mathbb{P}}$-local martingale, and a true martingale whenever the required integrability holds.
\end{enumerate}
When $B$ is the money-market account, $\widetilde{\mathbb{P}}$ is called a \emph{risk-neutral measure}.
\end{de}

\begin{thm}[First fundamental theorem of asset pricing, informal form]
Under the usual admissibility and regularity assumptions, the existence of an equivalent local martingale measure implies that the market admits no arbitrage. In a general semimartingale model, the no-free-lunch-with-vanishing-risk condition is equivalent to the existence of an equivalent local martingale measure.
\end{thm}

\subsection{One-period model}

Consider a market with $n$ possible terminal states and $m$ traded assets. Let
\[
p=(p_1,\ldots,p_m)^\top\in\mathbb{R}^m
\]
be the vector of time-zero prices, and let $X\in\mathbb{R}^{n\times m}$ be the payoff matrix, where $X_{ji}$ is the payoff of asset $i$ in state $j$. A portfolio $a\in\mathbb{R}^m$ costs $p^\top a$ at time zero and has terminal payoff vector $Xa$.

\begin{de}
A portfolio $a\in\mathbb{R}^m$ is an arbitrage if
\[
p^\top a\leq 0,
\qquad
Xa\geq 0,
\]
and at least one of these inequalities is strict. Inequalities between vectors are understood componentwise.
\end{de}

\begin{thm}[Finite-state fundamental theorem]
The one-period market is arbitrage-free if and only if there exists a strictly positive state-price vector $y\in\mathbb{R}_{++}^n$ such that
\[
p=X^\top y.
\]
\end{thm}

\begin{proof}
This is a finite-dimensional separating-hyperplane result, equivalently a form of the Farkas--Stiemke lemma. If $p=X^\top y$ for some $y>0$ and $a$ were an arbitrage, then
\[
p^\top a=y^\top Xa.
\]
The right-hand side is nonnegative, and it is strictly positive whenever $Xa$ is nonzero and nonnegative. This contradicts the arbitrage inequalities. The converse follows from strict separation of the attainable payoff cone from the nonnegative orthant.
\end{proof}

Suppose the market contains a risk-free asset with time-zero price $1$ and terminal payoff $1+r$, where $r>-1$. If $p=X^\top y$, then the pricing equation for the risk-free asset gives
\[
1=(1+r)\sum_{j=1}^n y_j.
\]
Hence
\[
D:=\sum_{j=1}^n y_j=\frac{1}{1+r},
\qquad
q_j:=\frac{y_j}{D},
\]
defines a probability vector $q=(q_1,\ldots,q_n)$ with strictly positive components. For each traded asset,
\[
p_i=D\sum_{j=1}^n q_jX_{ji}
   =\frac{1}{1+r}\,\widetilde{\mathbb{E}}[X_i].
\]
Thus $q$ is the one-period risk-neutral probability measure.

If a new claim has payoff vector $v\in\mathbb{R}^n$, every state-price vector consistent with traded prices assigns it the value
\[
\pi(v)=y^\top v
      =D\,\widetilde{\mathbb{E}}[v].
\]
When the state-price vector is not unique, a nonreplicable claim may have a range of arbitrage-free prices.

\begin{de}
The one-period market is \emph{complete} if every payoff vector $v\in\mathbb{R}^n$ can be replicated; that is, for every $v$ there exists $a\in\mathbb{R}^m$ such that
\[
Xa=v.
\]
Equivalently,
\[
\operatorname{rank}(X)=n.
\]
\end{de}

Completeness requires $m\geq n$, although some assets may be redundant. If $m<n$, then $\operatorname{rank}(X)<n$ and the market is necessarily incomplete. The inequality $m>n$ does not itself create arbitrage: it merely permits redundancy, provided prices are mutually consistent.

\begin{thm}[Second fundamental theorem in the finite-state model]
Assume the market is arbitrage-free and contains the risk-free asset described above. The market is complete if and only if the equivalent risk-neutral probability measure is unique.
\end{thm}

\subsection{Diffusive market model}

Let $(\Omega,\mathcal{F},(\mathcal{F}_t)_{t\geq 0},\mathbb{P})$ support a $d$-dimensional Brownian motion
\[
W_t=(W_t^{(1)},\ldots,W_t^{(d)})^\top.
\]
The money-market account satisfies
\[
dB_t=r_tB_t\,dt,
\qquad
B_0=1,
\]
so
\[
B_t=\exp\!\left(\int_0^t r_s\,ds\right),
\qquad
D_t:=B_t^{-1}=\exp\!\left(-\int_0^t r_s\,ds\right).
\]
The $m$ risky assets satisfy
\[
\frac{dS_t^{(i)}}{S_t^{(i)}}
 =\alpha_t^{(i)}\,dt
  +\sum_{j=1}^d\sigma_t^{(ij)}\,dW_t^{(j)},
\qquad i=1,\ldots,m.
\]
In vector form,
\[
\frac{dS_t}{S_t}=\alpha_t\,dt+\sigma_t\,dW_t,
\]
where the division is componentwise, $\alpha_t\in\mathbb{R}^m$, and $\sigma_t\in\mathbb{R}^{m\times d}$.

For asset $i$, the instantaneous variance is
\[
\operatorname{Var}_t\!\left(\frac{dS_t^{(i)}}{S_t^{(i)}}\right)
 =\sum_{k=1}^d\bigl(\sigma_t^{(ik)}\bigr)^2dt,
\]
and, for assets $i$ and $\ell$,
\[
\operatorname{Cov}_t\!\left(
\frac{dS_t^{(i)}}{S_t^{(i)}},
\frac{dS_t^{(\ell)}}{S_t^{(\ell)}}
\right)
 =\sum_{k=1}^d\sigma_t^{(ik)}\sigma_t^{(\ell k)}dt.
\]
Consequently, whenever the denominators are nonzero, their instantaneous correlation is
\[
\rho_t^{i\ell}
=\frac{\sum_{k=1}^d\sigma_t^{(ik)}\sigma_t^{(\ell k)}}
{\sqrt{\sum_{k=1}^d(\sigma_t^{(ik)})^2}
 \sqrt{\sum_{k=1}^d(\sigma_t^{(\ell k)})^2}}.
\]

Define the scalar instantaneous volatility of asset $i$ by
\[
\bar\sigma_t^{(i)}
 =\left(\sum_{j=1}^d(\sigma_t^{(ij)})^2\right)^{1/2}.
\]
When $\bar\sigma_t^{(i)}>0$, the process
\[
B_t^{(i)}
 =\int_0^t\sum_{j=1}^d
   \frac{\sigma_s^{(ij)}}{\bar\sigma_s^{(i)}}\,dW_s^{(j)}
\]
is a Brownian motion by L\'evy's characterization, and
\[
\frac{dS_t^{(i)}}{S_t^{(i)}}
 =\alpha_t^{(i)}dt+\bar\sigma_t^{(i)}dB_t^{(i)}.
\]
The Brownian motions $B^{(i)}$ need not be independent.

\subsubsection*{Trading strategies}

Let $\Delta_t^{(i)}$ be the number of shares held in risky asset $i$, let $\beta_t$ be the number of units held in the money-market account, and let
\[
X_t=\beta_tB_t+\sum_{i=1}^m\Delta_t^{(i)}S_t^{(i)}
\]
be total wealth. If consumption occurs at rate $c_t$, the self-financing-with-consumption condition is
\[
dX_t=\beta_t\,dB_t+\sum_{i=1}^m\Delta_t^{(i)}\,dS_t^{(i)}-c_t\,dt.
\]
Equivalently,
\[
dX_t
=r_tX_t\,dt
+\sum_{i=1}^m\Delta_t^{(i)}
  \bigl(dS_t^{(i)}-r_tS_t^{(i)}dt\bigr)
-c_t\,dt.
\]
When $c\equiv0$, the strategy is self-financing, and the discounted wealth satisfies
\[
d(D_tX_t)
 =\sum_{i=1}^m\Delta_t^{(i)}\,d(D_tS_t^{(i)}).
\]

\subsubsection*{Risk-neutral measures and Girsanov's theorem}

Let $\theta_t\in\mathbb{R}^d$ be progressively measurable and suppose the stochastic exponential
\[
Z_t
=\exp\!\left(
 -\int_0^t\theta_s^\top dW_s
 -\frac12\int_0^t\|\theta_s\|^2ds
\right)
\]
is a true martingale on $[0,T]$; Novikov's condition is one sufficient criterion. Define
\[
\frac{d\widetilde{\mathbb{P}}}{d\mathbb{P}}
\bigg|_{\mathcal{F}_T}=Z_T.
\]
Then
\[
\widetilde W_t=W_t+\int_0^t\theta_s\,ds
\]
is a $d$-dimensional Brownian motion under $\widetilde{\mathbb{P}}$. Substitution gives
\[
\frac{dS_t}{S_t}
=\bigl(\alpha_t-\sigma_t\theta_t\bigr)dt
 +\sigma_t\,d\widetilde W_t.
\]
Therefore $\widetilde{\mathbb{P}}$ is risk-neutral precisely when
\[
\boxed{\alpha_t-r_t\mathbf{1}_m=\sigma_t\theta_t}
\]
almost everywhere. Under this condition,
\[
\frac{dS_t^{(i)}}{S_t^{(i)}}
=r_tdt+\sum_{j=1}^d\sigma_t^{(ij)}d\widetilde W_t^{(j)}.
\]

\begin{thm}[No arbitrage from an equivalent martingale measure]
Suppose $\widetilde{\mathbb{P}}$ is an equivalent martingale measure and admissible discounted wealth processes are true martingales, or are local martingales bounded from below. Then no admissible self-financing arbitrage exists.
\end{thm}

\begin{proof}
For an admissible self-financing strategy,
\[
d(D_tX_t)
 =\sum_{i=1}^m\Delta_t^{(i)}d(D_tS_t^{(i)}).
\]
Thus $DX$ is a $\widetilde{\mathbb{P}}$-local martingale. Under the stated admissibility condition it is a supermartingale, and it is a martingale under the stronger integrability assumption. If $X_0=0$ and $X_T\geq0$ almost surely, then
\[
0\geq\widetilde{\mathbb{E}}[D_TX_T]\geq0,
\]
so $D_TX_T=0$ almost surely. Equivalence of the measures then rules out a strictly positive terminal payoff with positive $\mathbb{P}$-probability.
\end{proof}

\subsubsection*{Completeness}

\begin{de}
The market is complete on $[0,T]$ if every suitably integrable $\mathcal{F}_T$-measurable claim $H$ can be replicated by an admissible self-financing strategy; that is, there exists a wealth process $X$ such that $X_T=H$.
\end{de}

Assume that $(\mathcal{F}_t)$ is the augmented Brownian filtration generated by $\widetilde W$. If $D_TH\in L^2(\widetilde{\mathbb{P}})$, the martingale representation theorem gives a predictable process $\Gamma_t\in\mathbb{R}^d$ such that
\[
M_t:=\widetilde{\mathbb{E}}[D_TH\mid\mathcal{F}_t]
=M_0+\int_0^t\Gamma_s^\top d\widetilde W_s.
\]
A self-financing strategy has discounted dynamics
\[
d(D_tX_t)
=\sum_{i=1}^m\sum_{j=1}^d
 \Delta_t^{(i)}D_tS_t^{(i)}\sigma_t^{(ij)}
 \,d\widetilde W_t^{(j)}.
\]
Replication therefore requires
\[
\Gamma_t
=\sigma_t^\top
\begin{pmatrix}
\Delta_t^{(1)}D_tS_t^{(1)}\\
\vdots\\
\Delta_t^{(m)}D_tS_t^{(m)}
\end{pmatrix}.
\]
This linear system can be solved for every $\Gamma_t\in\mathbb{R}^d$ whenever
\[
\operatorname{rank}(\sigma_t)=d
\]
almost everywhere, which in particular requires $m\geq d$.

\begin{thm}[Second fundamental theorem, diffusion form]
Under the standard regularity assumptions for a Brownian diffusion market, the market is complete if and only if the equivalent martingale measure is unique. In the model above, this corresponds to the risk-premium equation
\[
\alpha_t-r_t\mathbf{1}_m=\sigma_t\theta_t
\]
having a unique solution $\theta_t$, equivalently $\sigma_t$ having full column rank $d$ almost everywhere.
\end{thm}

\subsection{Assets with dividends}

\subsubsection*{Continuous dividend yield}

Suppose one stock pays dividends continuously at rate $q_tS_tdt$, where $q_t\geq0$ is adapted. Under the physical measure, write its ex-dividend dynamics as
\[
\frac{dS_t}{S_t}=(\mu_t-q_t)dt+\sigma_t\,dW_t.
\]
The total gain from holding one share is
\[
dG_t=dS_t+q_tS_tdt
     =\mu_tS_tdt+\sigma_tS_tdW_t.
\]
A portfolio holding $\Delta_t$ shares therefore satisfies
\[
dX_t
=r_tX_tdt+\Delta_tS_t(\mu_t-r_t)dt
 +\Delta_t\sigma_tS_tdW_t.
\]
With
\[
\theta_t=\frac{\mu_t-r_t}{\sigma_t},
\qquad
\widetilde W_t=W_t+\int_0^t\theta_sds,
\]
the risk-neutral dynamics become
\[
\boxed{\frac{dS_t}{S_t}=(r_t-q_t)dt+\sigma_t\,d\widetilde W_t.}
\]
Thus the ex-dividend discounted stock price need not be a martingale. Instead, the discounted gains process
\[
\frac{S_t}{B_t}+\int_0^t\frac{q_uS_u}{B_u}\,du
\]
is a $\widetilde{\mathbb{P}}$-martingale. Equivalently, when dividends are continuously reinvested, the total-return process
\[
\widehat S_t=S_t\exp\!\left(\int_0^tq_u\,du\right)
\]
has discounted value $D_t\widehat S_t$ equal to a local martingale, and a true martingale under suitable integrability.

For a claim $H$ payable at $T$, risk-neutral valuation remains
\[
V_t
=\frac{1}{D_t}\widetilde{\mathbb{E}}[D_TH\mid\mathcal{F}_t].
\]
The dividend yield changes the risk-neutral dynamics of the underlying, not the basic discounting formula.

\subsubsection*{Discrete proportional dividends}

Let
\[
0<t_1<\cdots<t_n<T
\]
be dividend dates. Suppose the stock pays the proportional dividend
\[
a_jS_{t_j-},
\qquad 0\leq a_j\leq1,
\]
at time $t_j$, where $a_j$ is $\mathcal{F}_{t_j-}$-measurable. The no-arbitrage ex-dividend jump condition is
\[
S_{t_j}=(1-a_j)S_{t_j-}.
\]
Between dividend dates, suppose
\[
\frac{dS_t}{S_t}=\mu_tdt+\sigma_tdW_t.
\]
A shareholder receives the dividend at the same instant that the quoted ex-dividend price falls, so a self-financing portfolio's total value does not jump merely because of the payout. The discounted total gains process, rather than the discounted ex-dividend stock price alone, is the martingale under the risk-neutral measure.

\subsection{Derivatives pricing: forwards and futures}

\begin{de}
A forward contract initiated at time $t$ on an asset with maturity $T$ and delivery price $K$ pays
\[
S_T-K
\]
to the long position at time $T$. The forward price $\operatorname{For}_S(t,T)$ is the value of $K$ that makes the contract's value zero at inception.
\end{de}

Let
\[
P(t,T)
=\widetilde{\mathbb{E}}\!\left[
\frac{D_T}{D_t}\,\bigg|\,\mathcal{F}_t
\right]
\]
be the time-$t$ price of a zero-coupon bond paying one unit at $T$.

\begin{thm}
Assume the underlying is a non-dividend-paying traded asset and zero-coupon bonds are tradable. Then
\[
\boxed{\operatorname{For}_S(t,T)=\frac{S_t}{P(t,T)}.}
\]
\end{thm}

\begin{proof}
The value at time $t$ of a forward with delivery price $K$ is
\[
\frac{1}{D_t}
\widetilde{\mathbb{E}}[D_T(S_T-K)\mid\mathcal{F}_t]
=S_t-KP(t,T),
\]
because $D_tS_t$ is a martingale. Setting this value equal to zero gives the formula. Equivalently, one can buy the asset and finance it by shorting $S_t/P(t,T)$ units of the $T$-bond.
\end{proof}

For a dividend-paying asset, $S_t$ in this formula must be replaced by the time-$t$ prepaid-forward price, namely the current asset price minus the present value of dividends paid before $T$ when those dividends are known.

\begin{de}
A futures contract is marked to market. Its quoted price $\operatorname{Fut}_S(t,T)$ is reset through gains and losses in a margin account. Over an interval $(u,v]$, the long position receives
\[
\operatorname{Fut}_S(v,T)-\operatorname{Fut}_S(u,T).
\]
At maturity,
\[
\operatorname{Fut}_S(T,T)=S_T.
\]
\end{de}

An increase in the futures quotation benefits the long position and harms the short position. Futures are standardized and exchange-traded; forwards are usually customized and traded over the counter. Over-the-counter contracts are still tradable, but they generally carry more counterparty-credit and liquidity risk.

\begin{thm}
Under continuous marking to market and the usual frictionless assumptions, the futures price is
\[
\boxed{\operatorname{Fut}_S(t,T)
=\widetilde{\mathbb{E}}[S_T\mid\mathcal{F}_t].}
\]
In particular, the futures-price process is a $\widetilde{\mathbb{P}}$-martingale.
\end{thm}

\begin{proof}
A futures contract has zero value immediately after each marking-to-market payment, and a position of $\Gamma_t$ contracts contributes the gain $\Gamma_t\,d\operatorname{Fut}_S(t,T)$ to wealth. For discounted wealth to remain a local martingale under the risk-neutral measure for arbitrary admissible $\Gamma$, the futures quotation must itself be a local martingale. The terminal condition then yields
\[
\operatorname{Fut}_S(t,T)
=\widetilde{\mathbb{E}}[\operatorname{Fut}_S(T,T)\mid\mathcal{F}_t]
=\widetilde{\mathbb{E}}[S_T\mid\mathcal{F}_t].
\]
\end{proof}

If interest rates are deterministic, then
\[
\operatorname{For}_S(t,T)
=\operatorname{Fut}_S(t,T).
\]
At time zero, with stochastic interest rates,
\[
\begin{aligned}
\operatorname{For}_S(0,T)-\operatorname{Fut}_S(0,T)
&=\frac{\widetilde{\mathbb{E}}[D_TS_T]}
        {\widetilde{\mathbb{E}}[D_T]}
  -\widetilde{\mathbb{E}}[S_T]\\
&=\frac{\widetilde{\operatorname{Cov}}(D_T,S_T)}{P(0,T)}.
\end{aligned}
\]
This covariance formula explains why stochastic rates can make forward and futures prices differ.

\subsection{Derivatives pricing: options}

\begin{de}
A European call gives its holder the right, but not the obligation, to buy the underlying asset for the \emph{strike price} $K$ at maturity $T$. Its payoff is
\[
(S_T-K)^+.
\]
A European put gives the right to sell at $K$ and has payoff
\[
(K-S_T)^+.
\]
\end{de}

For a European payoff $h(S_T)$, the arbitrage-free value in a complete market is
\[
V_t
=\frac{1}{D_t}
 \widetilde{\mathbb{E}}[D_Th(S_T)\mid\mathcal{F}_t].
\]
In an incomplete market this formula depends on the selected equivalent martingale measure unless the claim is replicable.

\subsection{The Black--Scholes formula}

Assume a non-dividend-paying stock with physical dynamics
\[
\frac{dS_t}{S_t}=\mu\,dt+\sigma\,dW_t,
\qquad \sigma>0,
\]
and a constant risk-free rate $r$. Set
\[
\theta=\frac{\mu-r}{\sigma},
\qquad
Z_t=\exp\!\left(-\theta W_t-\frac12\theta^2t\right).
\]
Under the measure $\widetilde{\mathbb{P}}$ defined by $Z_T$, the process
\[
\widetilde W_t=W_t+\theta t
\]
is Brownian motion and
\[
\frac{dS_t}{S_t}=r\,dt+\sigma\,d\widetilde W_t.
\]

\begin{thm}[Black--Scholes formula]
For $0\leq t<T$, the European call price is
\[
C_t=S_tN(d_1)-Ke^{-r(T-t)}N(d_2),
\]
where
\[
d_1
=\frac{\log(S_t/K)+(r+\frac12\sigma^2)(T-t)}
       {\sigma\sqrt{T-t}},
\qquad
d_2=d_1-\sigma\sqrt{T-t},
\]
and $N$ is the standard normal distribution function.
\end{thm}

\begin{proof}
Let $\tau=T-t$. Under $\widetilde{\mathbb{P}}$,
\[
S_T
=S_t\exp\!\left(
(r-\tfrac12\sigma^2)\tau
+\sigma(\widetilde W_T-\widetilde W_t)
\right).
\]
Decompose the payoff:
\[
C_t
=\widetilde{\mathbb{E}}\!\left[
 e^{-r\tau}S_T\mathbf{1}_{\{S_T>K\}}
 \mid\mathcal{F}_t
\right]
-Ke^{-r\tau}
 \widetilde{\mathbb{P}}(S_T>K\mid\mathcal{F}_t).
\]
For the first term, introduce the stock-numeraire measure $\mathbb{P}^S$ by the density process
\[
L_u^S=\frac{e^{-ru}S_u}{S_0},
\qquad 0\leq u\leq T.
\]
Then Bayes' formula gives
\[
\widetilde{\mathbb{E}}\!\left[
 e^{-r\tau}S_T\mathbf{1}_{\{S_T>K\}}
 \mid\mathcal{F}_t
\right]
=S_t\mathbb{P}^S(S_T>K\mid\mathcal{F}_t).
\]
Under $\mathbb{P}^S$,
\[
W_t^S=\widetilde W_t-\sigma t
\]
is Brownian motion, and the conditional lognormal calculations yield
\[
\mathbb{P}^S(S_T>K\mid\mathcal{F}_t)=N(d_1),
\qquad
\widetilde{\mathbb{P}}(S_T>K\mid\mathcal{F}_t)=N(d_2).
\]
Substitution proves the formula.
\end{proof}

\begin{thm}[Put--call parity]
For a non-dividend-paying stock and constant interest rate $r$,
\[
C_t-P_t=S_t-Ke^{-r(T-t)}.
\]
\end{thm}

\begin{proof}
The identity
\[
(S_T-K)^+-(K-S_T)^+=S_T-K
\]
and risk-neutral valuation imply
\[
C_t-P_t
=e^{-r(T-t)}\widetilde{\mathbb{E}}[S_T-K\mid\mathcal{F}_t]
=S_t-Ke^{-r(T-t)}.
\]
\end{proof}

\begin{thm}[Black--Scholes PDE]
If $C_t=c(t,S_t)$ and $c$ is sufficiently smooth, then
\[
c_t(t,x)+rx c_x(t,x)
+\frac12\sigma^2x^2c_{xx}(t,x)-rc(t,x)=0,
\]
with terminal condition
\[
c(T,x)=(x-K)^+.
\]
\end{thm}

\begin{proof}
Under the risk-neutral measure,
\[
dS_t=rS_tdt+\sigma S_td\widetilde W_t.
\]
It\^o's formula gives
\[
dc(t,S_t)
=\left(c_t+rS_tc_x+\frac12\sigma^2S_t^2c_{xx}\right)dt
 +\sigma S_tc_xd\widetilde W_t.
\]
Therefore
\[
\begin{aligned}
d(e^{-rt}c(t,S_t))
&=e^{-rt}
 \left(c_t+rS_tc_x+\frac12\sigma^2S_t^2c_{xx}-rc\right)dt\\
&\quad+e^{-rt}\sigma S_tc_xd\widetilde W_t.
\end{aligned}
\]
The discounted option value is a martingale, so its drift must vanish.
\end{proof}

A replicating portfolio starts with $X_0=c(0,S_0)$ and holds
\[
\Delta_t=c_x(t,S_t)
\]
shares of stock. Indeed,
\[
d(e^{-rt}X_t)=e^{-rt}\sigma S_t\Delta_t\,d\widetilde W_t
\]
after matching the option's discounted diffusion term. For a Black--Scholes call,
\[
\Delta_t=N(d_1).
\]

The principal option sensitivities are
\[
\Delta=c_x,
\qquad
\Gamma=c_{xx},
\qquad
\Theta=c_t,
\qquad
\rho=\frac{\partial c}{\partial r},
\qquad
\mathcal{V}=\frac{\partial c}{\partial\sigma},
\]
where $\mathcal{V}$ is the vega.

\subsection{Limitations of the Black--Scholes model}

Market prices are commonly summarized by \emph{implied volatility}. For fixed $(t,S_t,K,T,r)$, the Black--Scholes call price is continuous and strictly increasing in $\sigma>0$. Hence a unique finite positive implied volatility exists when the observed call price lies strictly inside the model's no-arbitrage range
\[
(S_t-Ke^{-r(T-t)})^+<C_t^{\mathrm{mkt}}<S_t.
\]
At a boundary, the corresponding limiting volatility may be zero or infinite; outside the bounds, no Black--Scholes implied volatility exists.

In the Black--Scholes model, one constant volatility must fit all strikes and maturities. Market implied volatilities instead vary with both strike and maturity, producing smiles or skews. Models designed to capture these patterns include local-volatility, stochastic-volatility, and jump models; interest-rate term-structure models address the evolution of rates rather than, by themselves, the equity implied-volatility surface.

\begin{de}
A local-volatility model has risk-neutral dynamics
\[
dS_t=rS_tdt+\sigma_{\mathrm{loc}}(t,S_t)S_t\,d\widetilde W_t.
\]
A common constant-elasticity-of-variance specification is
\[
dS_t=rS_tdt+\sigma_0S_t^{\gamma}\,d\widetilde W_t.
\]
When $\gamma<1$, the percentage volatility $\sigma_0S_t^{\gamma-1}$ rises as the stock price falls, producing a leverage effect.
\end{de}

\begin{de}
A stochastic-volatility model has dynamics of the form
\[
\begin{aligned}
dS_t&=rS_tdt+\sqrt{v_t}\,S_t\,d\widetilde W_t^{(1)},\\
dv_t&=a(v_t)dt+b(v_t)d\widetilde W_t^{(2)},\\
d\widetilde W_t^{(1)}d\widetilde W_t^{(2)}&=\rho\,dt.
\end{aligned}
\]
The Heston model specifies
\[
dv_t=\kappa(\theta-v_t)dt+\xi\sqrt{v_t}\,d\widetilde W_t^{(2)}.
\]
A negative correlation $\rho$ generates a leverage effect and typically an equity-volatility skew.
\end{de}

\section{Change of numeraire}

\subsection{Motivation and general properties}

A martingale measure is always associated with a chosen numeraire. If the money-market account $B$ is used as numeraire, the corresponding risk-neutral measure $Q$ makes every discounted, non-dividend-paying traded-asset price $S_t/B_t$ a local martingale. In many applications, changing from $B$ to another strictly positive traded asset can simplify pricing calculations substantially.

Consider an arbitrage-free market with asset prices
\[
S^0,S^1,\ldots,S^n,
\]
where the candidate numeraires are strictly positive and pay no dividends. When a Brownian representation is needed, suppose that under a reference measure the dynamics are
\[
\frac{dS_t^i}{S_t^i}
=\mu_t^i\,dt+(\sigma_t^i)^\top dW_t,
\qquad i=0,\ldots,n,
\]
where $W$ is multidimensional and the coefficient processes are adapted.

A numeraire must be more than a positive deflator: it must be the value process of a strictly positive traded asset or self-financing portfolio. A nonfinancial index, a futures quotation, or the ex-dividend price of a dividend-paying asset cannot automatically be used as a numeraire without a separate argument.

Suppose $Q^0$ is the martingale measure associated with $S^0$. Thus, for every traded asset price $\Pi$,
\[
\frac{\Pi_t}{S_t^0}
\]
is a $Q^0$-local martingale, and a true martingale under the integrability conditions used below. We seek the martingale measure associated with $S^1$ while preserving the same arbitrage-free price system.

\begin{prop}[Change-of-numeraire theorem]
Assume that $S^1/S^0$ is a strictly positive true $Q^0$-martingale on $[0,T]$. Define
\[
L_t^{01}
=\frac{S_t^1/S_t^0}{S_0^1/S_0^0},
\qquad 0\leq t\leq T,
\]
and define $Q^1$ on $\mathcal{F}_T$ by
\[
\frac{dQ^1}{dQ^0}\bigg|_{\mathcal{F}_T}=L_T^{01}.
\]
Then $Q^1$ is the martingale measure associated with the numeraire $S^1$, and $Q^0$ and $Q^1$ generate the same price system.
\end{prop}

\begin{proof}
Let $\Pi$ be an arbitrage-free price process such that the required conditional expectations are finite. For $s\leq t$, Bayes' formula gives
\[
\begin{aligned}
E^{1}\!\left[
 \frac{\Pi_t}{S_t^1}\bigg|\mathcal{F}_s
\right]
&=\frac{1}{L_s^{01}}
 E^{0}\!\left[
 L_t^{01}\frac{\Pi_t}{S_t^1}
 \bigg|\mathcal{F}_s
 \right]\\
&=\frac{1}{L_s^{01}}
 \frac{S_0^0}{S_0^1}
 E^{0}\!\left[
 \frac{\Pi_t}{S_t^0}
 \bigg|\mathcal{F}_s
 \right]\\
&=\frac{1}{L_s^{01}}
 \frac{S_0^0}{S_0^1}
 \frac{\Pi_s}{S_s^0}
=\frac{\Pi_s}{S_s^1}.
\end{aligned}
\]
Thus $\Pi/S^1$ is a $Q^1$-martingale.

For a claim $X$ payable at $T$,
\[
\begin{aligned}
S_t^1E^1\!\left[\frac{X}{S_T^1}\bigg|\mathcal{F}_t\right]
&=S_t^1\frac{1}{L_t^{01}}
 E^0\!\left[L_T^{01}\frac{X}{S_T^1}\bigg|\mathcal{F}_t\right]\\
&=S_t^0E^0\!\left[\frac{X}{S_T^0}\bigg|\mathcal{F}_t\right],
\end{aligned}
\]
so the two numeraires generate the same prices.
\end{proof}

\begin{cor}
Let $Q$ be the risk-neutral measure associated with the money-market account $B$, and let $S$ be a strictly positive non-dividend-paying traded asset. The $S$-numeraire measure $Q^S$ has density process
\[
L_t^S
=\frac{dQ^S}{dQ}\bigg|_{\mathcal{F}_t}
=\frac{S_t/B_t}{S_0/B_0}.
\]
If $B_0=1$, this becomes
\[
L_t^S=\frac{S_t}{B_tS_0}.
\]
\end{cor}

\subsubsection*{Girsanov kernels}

Suppose that, under $Q^0$, the density process satisfies
\[
\frac{dL_t^{01}}{L_t^{01}}
=(\sigma_t^1-\sigma_t^0)^\top dW_t^0.
\]
Then
\[
L_t^{01}
=\mathcal{E}\!\left(
 \int_0^t(\sigma_s^1-\sigma_s^0)^\top dW_s^0
\right),
\]
and the $Q^1$-Brownian motion is
\[
W_t^1
=W_t^0-
 \int_0^t(\sigma_s^1-\sigma_s^0)\,ds.
\]
Thus, with the density convention above, the Girsanov kernel for the change from $Q^0$ to $Q^1$ is the volatility difference
\[
\varphi_t^{01}=\sigma_t^1-\sigma_t^0.
\]

In particular, suppose that under $Q$,
\[
\frac{dS_t}{S_t}=r_tdt+\sigma_t^\top dW_t^Q.
\]
Then
\[
\frac{dL_t^S}{L_t^S}=\sigma_t^\top dW_t^Q,
\]
and
\[
W_t^S=W_t^Q-\int_0^t\sigma_sds
\]
is Brownian motion under $Q^S$.

\subsection{Forward measures}

Fix a maturity $T$, and let $p(t,T)$ denote the time-$t$ price of a zero-coupon bond paying one unit at time $T$. The bond is strictly positive for $t\leq T$ and satisfies
\[
p(T,T)=1.
\]

\begin{de}
The $T$-forward measure $Q^T$ is the martingale measure associated with the numeraire $p(t,T)$.
\end{de}

\begin{prop}
Let $Q$ be the risk-neutral measure associated with $B$. The density process of $Q^T$ relative to $Q$ is
\[
L_t^T
=\frac{dQ^T}{dQ}\bigg|_{\mathcal{F}_t}
=\frac{p(t,T)/B_t}{p(0,T)/B_0}.
\]
If $B_0=1$, then
\[
L_t^T=\frac{p(t,T)}{B_tp(0,T)}.
\]

Suppose the $Q$-dynamics of the $T$-bond are
\[
\frac{dp(t,T)}{p(t,T)}
=r_tdt+v(t,T)^\top dW_t^Q,
\]
where $v(t,T)\in\mathbb{R}^d$. Then
\[
\frac{dL_t^T}{L_t^T}
=v(t,T)^\top dW_t^Q,
\]
and
\[
W_t^T
=W_t^Q-\int_0^tv(s,T)ds
\]
is Brownian motion under $Q^T$.
\end{prop}

\begin{proof}
The density formula is the change-of-numeraire theorem with $S=p(\cdot,T)$. Solving the bond SDE gives
\[
\begin{aligned}
p(t,T)
&=p(0,T)
 \exp\!\left(
 \int_0^t\left(r_s-\frac12\|v(s,T)\|^2\right)ds
 +\int_0^tv(s,T)^\top dW_s^Q
 \right).
\end{aligned}
\]
Since
\[
B_t=B_0\exp\!\left(\int_0^tr_sds\right),
\]
we obtain
\[
L_t^T
=\exp\!\left(
 -\frac12\int_0^t\|v(s,T)\|^2ds
 +\int_0^tv(s,T)^\top dW_s^Q
\right),
\]
which has the claimed differential.
\end{proof}

\begin{prop}[Pricing under the forward measure]
For any integrable claim $X$ payable at time $T$,
\[
\boxed{
\Pi_t[X]
=p(t,T)E^T[X\mid\mathcal{F}_t].}
\]
\end{prop}

\begin{proof}
Under the $T$-forward measure, every traded price divided by $p(\cdot,T)$ is a martingale. Since the claim has value $X$ at $T$ and $p(T,T)=1$,
\[
\frac{\Pi_t[X]}{p(t,T)}
=E^T\!\left[
 \frac{X}{p(T,T)}\bigg|\mathcal{F}_t
\right]
=E^T[X\mid\mathcal{F}_t].
\]
\end{proof}

\subsection{The numeraire portfolio}

Assume a money-market account $B$ and a fixed risk-neutral measure $Q$ have been chosen. Suppose an equivalent probability measure $Q^*$ is given, with density process
\[
L_t^*
=\frac{dQ^*}{dQ}\bigg|_{\mathcal{F}_t}.
\]
Since $L^*$ is a strictly positive true $Q$-martingale with $L_0^*=1$, define
\[
N_t=\frac{B_t}{B_0}L_t^*.
\]
Then $N_0=1$ and
\[
\frac{N_t}{B_t}=\frac{L_t^*}{B_0}
\]
is a $Q$-martingale. Therefore $N$ is a strictly positive arbitrage-free price process in the price system generated by $Q$.

\begin{prop}
If the process $N$ defined above is the value process of a strictly positive traded self-financing portfolio, then it is a legitimate numeraire and
\[
Q^*=Q^N.
\]
In a complete market, every suitably integrable arbitrage-free price process is attainable, so this condition is automatic. In an incomplete market, attainability of $N$ must be checked separately.
\end{prop}

\begin{proof}
The density of the $N$-numeraire measure relative to $Q$ is
\[
\frac{N_t/B_t}{N_0/B_0}
=\frac{(L_t^*/B_0)}{1/B_0}
=L_t^*.
\]
Hence the two measures agree on every $\mathcal{F}_t$. The fact that $L^*$ is a true martingale follows from its definition as a density process; it is not a consequence merely of working on a finite time horizon.
\end{proof}

\subsection{Application: a general option-pricing formula}

Let $S$ be a strictly positive, non-dividend-paying traded asset, let $p(t,T)$ be the $T$-bond price, and consider a European call with payoff
\[
(S_T-K)^+.
\]

\begin{prop}[General option-pricing formula]
The time-$t$ call price is
\[
\boxed{
c_t
=S_tQ^S(S_T\geq K\mid\mathcal{F}_t)
-Kp(t,T)Q^T(S_T\geq K\mid\mathcal{F}_t).}
\]
\end{prop}

\begin{proof}
Write
\[
(S_T-K)^+
=S_T\mathbf{1}_{\{S_T\geq K\}}
-K\mathbf{1}_{\{S_T\geq K\}}.
\]
Use $S$ as numeraire for the first claim and the $T$-bond as numeraire for the second. The change-of-numeraire pricing formula gives
\[
\begin{aligned}
\Pi_t\!\left[S_T\mathbf{1}_{\{S_T\geq K\}}\right]
&=S_tQ^S(S_T\geq K\mid\mathcal{F}_t),\\
\Pi_t\!\left[\mathbf{1}_{\{S_T\geq K\}}\right]
&=p(t,T)Q^T(S_T\geq K\mid\mathcal{F}_t).
\end{aligned}
\]
Subtracting the two values proves the result.
\end{proof}

Suppose the relative price
\[
Z_{S,T}(u)=\frac{S_u}{p(u,T)}
\]
has deterministic volatility $\sigma_{S,T}(u)\in\mathbb{R}^d$ on $[t,T]$. Define
\[
\Sigma_{S,T}^2(t,T)
=\int_t^T\|\sigma_{S,T}(u)\|^2du.
\]

\begin{thm}[Geman--El Karoui--Rochet formula]
Under the preceding deterministic-volatility assumption,
\[
c_t=S_tN(d_1)-Kp(t,T)N(d_2),
\]
where
\[
\begin{aligned}
d_2
&=\frac{
 \log\!\left(\dfrac{S_t}{Kp(t,T)}\right)
 -\frac12\Sigma_{S,T}^2(t,T)
}{\sqrt{\Sigma_{S,T}^2(t,T)}},\\[4pt]
d_1
&=d_2+\sqrt{\Sigma_{S,T}^2(t,T)}.
\end{aligned}
\]
\end{thm}

\begin{proof}
Under the $T$-forward measure, the relative price is a martingale and satisfies
\[
\frac{dZ_{S,T}(u)}{Z_{S,T}(u)}
=\sigma_{S,T}(u)^\top dW_u^T.
\]
Hence
\[
Z_{S,T}(T)
=\frac{S_t}{p(t,T)}
 \exp\!\left(
 -\frac12\Sigma_{S,T}^2(t,T)
 +\int_t^T\sigma_{S,T}(u)^\top dW_u^T
 \right).
\]
Since $p(T,T)=1$, the event $S_T\geq K$ is the event
$Z_{S,T}(T)\geq K$. The stochastic integral is Gaussian with variance
$\Sigma_{S,T}^2(t,T)$, so
\[
Q^T(S_T\geq K\mid\mathcal{F}_t)=N(d_2).
\]

Under the stock measure, the reciprocal process
\[
Y_{S,T}(u)=\frac{p(u,T)}{S_u}=\frac1{Z_{S,T}(u)}
\]
is a martingale with volatility $-\sigma_{S,T}(u)$. Thus
\[
Y_{S,T}(T)
=\frac{p(t,T)}{S_t}
 \exp\!\left(
 -\frac12\Sigma_{S,T}^2(t,T)
 -\int_t^T\sigma_{S,T}(u)^\top dW_u^S
 \right).
\]
The event $S_T\geq K$ is equivalent to $Y_{S,T}(T)\leq1/K$, which yields
\[
Q^S(S_T\geq K\mid\mathcal{F}_t)=N(d_1).
\]
Substitution into the general option-pricing formula proves the result.
\end{proof}



    \section{Exotic options}

European calls and puts are often called \emph{vanilla} options because their payoffs depend only on the terminal value of the underlying asset. An option whose payoff or exercise rule depends on the path of the underlying asset is called \emph{path-dependent}; many such contracts are also called \emph{exotic options}.

Throughout most of this section, the underlying stock follows the Black--Scholes dynamics under the risk-neutral measure $\widetilde{\mathbb{P}}$:
\[
dS_t=rS_t\,dt+\sigma S_t\,d\widetilde W_t,
\qquad r>0,
\quad \sigma>0.
\]

\subsection{Barrier options}

A barrier option is activated or extinguished when the underlying reaches a specified level. A knock-out option becomes worthless after the barrier is reached, whereas a knock-in option becomes active only after the barrier is reached. The qualifiers \emph{up} and \emph{down} indicate whether the barrier lies above or below the initial stock price.

The running maximum of a drifted Brownian motion is useful for pricing an up-and-out option.

\begin{thm}
Let
\[
\widehat W_t=\alpha t+\widetilde W_t,
\qquad
\widehat M_T=\max_{0\leq u\leq T}\widehat W_u.
\]
Under $\widetilde{\mathbb{P}}$, the joint density of $(\widehat M_T,\widehat W_T)$ is
\[
\widetilde f_{\widehat M_T,\widehat W_T}(m,w)
=\frac{2(2m-w)}{T\sqrt{2\pi T}}
 \exp\!\left(
 \alpha w-\frac12\alpha^2T-
 \frac{(2m-w)^2}{2T}
 \right)
\]
for $m\geq0$ and $w\leq m$, and is zero elsewhere.
\end{thm}

\begin{proof}
Define
\[
\widehat Z_t
=\exp\!\left(-\alpha\widetilde W_t-\frac12\alpha^2t\right)
=\exp\!\left(-\alpha\widehat W_t+\frac12\alpha^2t\right)
\]
and let $\widehat{\mathbb{P}}$ be given by
\[
\frac{d\widehat{\mathbb{P}}}{d\widetilde{\mathbb{P}}}
\bigg|_{\mathcal{F}_T}=\widehat Z_T.
\]
Girsanov's theorem implies that $\widehat W$ is a standard Brownian motion under $\widehat{\mathbb{P}}$. By the reflection principle, its maximum and terminal value have joint density
\[
\widehat f_{\widehat M_T,\widehat W_T}(m,w)
=\frac{2(2m-w)}{T\sqrt{2\pi T}}
 \exp\!\left(-\frac{(2m-w)^2}{2T}\right)
\]
after restricting to $m\geq0$ and $w\leq m$. Since
\[
\frac{d\widetilde{\mathbb{P}}}{d\widehat{\mathbb{P}}}
\bigg|_{\mathcal{F}_T}
=\widehat Z_T^{-1}
=\exp\!\left(\alpha\widehat W_T-\frac12\alpha^2T\right),
\]
the change-of-measure formula multiplies the preceding density by
$\exp(\alpha w-\frac12\alpha^2T)$, giving the stated result.
\end{proof}

\begin{cor}
For $m\geq0$,
\[
\widetilde{\mathbb{P}}(\widehat M_T\leq m)
=N\!\left(\frac{m-\alpha T}{\sqrt T}\right)
-e^{2\alpha m}
 N\!\left(\frac{-m-\alpha T}{\sqrt T}\right).
\]
Its density is
\[
\widetilde f_{\widehat M_T}(m)
=\frac{2}{\sqrt{2\pi T}}
 \exp\!\left(-\frac{(m-\alpha T)^2}{2T}\right)
-2\alpha e^{2\alpha m}
 N\!\left(\frac{-m-\alpha T}{\sqrt T}\right),
\qquad m\geq0.
\]
\end{cor}

Now consider a continuously monitored up-and-out call with strike $K$, barrier $B$, and maturity $T$, where $0<S_0<B$. Under the risk-neutral measure,
\[
S_u
=S_0\exp\!\left(
(r-\tfrac12\sigma^2)u+\sigma\widetilde W_u
\right)
=S_0e^{\sigma\widehat W_u},
\]
where
\[
\alpha=\frac{r-\frac12\sigma^2}{\sigma}.
\]
Set
\[
k=\frac{1}{\sigma}\log\frac{K}{S_0},
\qquad
b=\frac{1}{\sigma}\log\frac{B}{S_0}.
\]
The payoff is
\[
V_T=(S_T-K)^+
\mathbf{1}_{\{\max_{0\leq u\leq T}S_u<B\}}
=\bigl(S_0e^{\sigma\widehat W_T}-K\bigr)
 \mathbf{1}_{\{\widehat W_T\geq k,\ \widehat M_T<b\}}.
\]
If $K\geq B$, the continuously monitored up-and-out call is worthless. If $K<B$, its time-zero value can be written as
\[
\begin{aligned}
V_0
&=e^{-rT}
 \int_k^b\int_{\max\{w,0\}}^b
 \bigl(S_0e^{\sigma w}-K\bigr)
 \widetilde f_{\widehat M_T,\widehat W_T}(m,w)
 \,dm\,dw.
\end{aligned}
\]
This integral can be reduced to normal distribution functions by completing squares.

\subsubsection*{PDE characterization}

Let
\[
\rho_B=\inf\{u\geq0:S_u\geq B\}
\]
be the first barrier-hitting time. Before knock-out, the discounted option value stopped at $\rho_B$ is a martingale.

\begin{thm}
Let $v(t,x)$ be the value at time $t$ of the up-and-out call conditional on $S_t=x$ and on no previous knock-out. In the continuation region
\[
0\leq t<T,
\qquad 0<x<B,
\]
the function $v$ satisfies
\[
v_t+rxv_x+\frac12\sigma^2x^2v_{xx}-rv=0.
\]
The boundary and terminal conditions are
\[
\begin{aligned}
v(t,0)&=0, &&0\leq t\leq T,\\
v(t,B)&=0, &&0\leq t<T,\\
v(T,x)&=(x-K)^+, &&0\leq x<B.
\end{aligned}
\]
At the corner $(T,B)$, continuous monitoring makes the barrier condition authoritative: a path that reaches $B$ at maturity is knocked out.
\end{thm}

\begin{proof}
For $t<\rho_B$, It\^o's formula gives
\[
\begin{aligned}
d\bigl(e^{-rt}v(t,S_t)\bigr)
&=e^{-rt}
 \left(v_t+rS_tv_x+\frac12\sigma^2S_t^2v_{xx}-rv\right)dt\\
&\quad+e^{-rt}\sigma S_tv_x\,d\widetilde W_t.
\end{aligned}
\]
The stopped discounted value is a martingale, so the drift vanishes in the interior. The remaining conditions follow from the payoff and the knock-out rule.
\end{proof}

A replicating strategy before knock-out holds
\[
\Delta_t=v_x(t,S_t)
\]
shares of stock. The hedge is stopped at $T\wedge\rho_B$.

\subsection{Lookback options}

Let
\[
Y_t=\max_{0\leq u\leq t}S_u
\]
be the running maximum. The payoff
\[
Y_T-S_T
\]
is a floating-strike lookback \emph{put}; a floating-strike lookback call is based on the running minimum.

The risk-neutral value is
\[
V_t
=\widetilde{\mathbb{E}}\!\left[
 e^{-r(T-t)}(Y_T-S_T)
 \mid\mathcal{F}_t
\right].
\]
Because $(S,Y)$ is Markov, there is a function $v$ such that
\[
V_t=v(t,S_t,Y_t).
\]

\begin{thm}
For $0\leq t<T$ and $0<x<y$, the floating-strike lookback-put value satisfies
\[
v_t+rxv_x+\frac12\sigma^2x^2v_{xx}-rv=0.
\]
The boundary and terminal conditions are
\[
\begin{aligned}
v(t,0,y)&=e^{-r(T-t)}y,
    &&0\leq t\leq T,\ y\geq0,\\
v_y(t,y,y)&=0,
    &&0\leq t<T,\ y>0,\\
v(T,x,y)&=y-x,
    &&0\leq x\leq y.
\end{aligned}
\]
\end{thm}

\begin{proof}
The running maximum $Y$ is continuous and nondecreasing, hence has finite variation and zero quadratic variation. Moreover, its Stieltjes measure $dY_t$ is supported on the set
\[
\{t:S_t=Y_t\};
\]
that is, $Y$ can increase only when the stock is at its current maximum.

It\^o's formula yields
\[
\begin{aligned}
d\bigl(e^{-rt}v(t,S_t,Y_t)\bigr)
&=e^{-rt}
 \left(v_t+rS_tv_x+\frac12\sigma^2S_t^2v_{xx}-rv\right)dt\\
&\quad+e^{-rt}\sigma S_tv_x\,d\widetilde W_t
 +e^{-rt}v_y(t,S_t,Y_t)\,dY_t.
\end{aligned}
\]
The $dt$ drift vanishes in the interior. Since $dY_t$ is carried by $S_t=Y_t$, the finite-variation term disappears precisely when
\[
v_y(t,y,y)=0.
\]
The terminal condition is the payoff. If $S_t=0$ and the zero boundary is treated as absorbing, then $S_u=0$ and $Y_u=y$ for $u\geq t$, which gives $v(t,0,y)=e^{-r(T-t)}y$.
\end{proof}

Once the PDE and boundary conditions hold,
\[
d\bigl(e^{-rt}v(t,S_t,Y_t)\bigr)
=e^{-rt}\sigma S_tv_x(t,S_t,Y_t)\,d\widetilde W_t,
\]
so the delta hedge is
\[
\Delta_t=v_x(t,S_t,Y_t).
\]

\subsection{Asian options}

An Asian option depends on an average of the underlying price. A continuously sampled arithmetic average is
\[
\frac{1}{T}\int_0^TS_u\,du,
\]
whereas a discretely sampled average is
\[
\frac{1}{m}\sum_{j=1}^mS_{t_j},
\qquad
0<t_1<\cdots<t_m=T.
\]

Consider a fixed-strike continuously sampled Asian call with payoff
\[
V_T=\left(\frac1T\int_0^TS_u\,du-K\right)^+.
\]
Define the accumulated integral
\[
Y_t=\int_0^tS_u\,du,
\qquad
dY_t=S_t\,dt.
\]
The augmented state $(S_t,Y_t)$ is Markov, and
\[
V_t=v(t,S_t,Y_t)
\]
for some deterministic function $v$.

\begin{thm}
The Asian-call price function satisfies
\[
v_t+rxv_x+xv_y+\frac12\sigma^2x^2v_{xx}-rv=0,
\qquad
0\leq t<T,
\quad x>0.
\]
Its terminal condition is
\[
v(T,x,y)=\left(\frac{y}{T}-K\right)^+.
\]
A natural boundary condition at $x=0$ is
\[
v(t,0,y)
=e^{-r(T-t)}\left(\frac{y}{T}-K\right)^+.
\]
If the PDE is extended mathematically to $y\in\mathbb{R}$, an additional natural condition is
\[
\lim_{y\downarrow-\infty}v(t,x,y)=0.
\]
\end{thm}

\begin{proof}
Because $dY_t=S_tdt$, we have $dS_t\,dY_t=dY_t\,dY_t=0$. It\^o's formula gives
\[
\begin{aligned}
d\bigl(e^{-rt}v(t,S_t,Y_t)\bigr)
&=e^{-rt}
 \left(v_t+rS_tv_x+S_tv_y
 +\frac12\sigma^2S_t^2v_{xx}-rv\right)dt\\
&\quad+e^{-rt}\sigma S_tv_x\,d\widetilde W_t.
\end{aligned}
\]
The discounted price is a martingale, so the drift vanishes. The terminal and boundary conditions follow directly from the payoff and from the fact that $S$ remains at zero once the absorbing boundary is reached.
\end{proof}

The replicating delta is again
\[
\Delta_t=v_x(t,S_t,Y_t).
\]

\subsection{American options: perpetual options}

The perpetual American put has no fixed maturity. If exercised at a stopping time $\tau$, it pays $(K-S_\tau)^+$. Let $\mathcal{T}$ denote the set of stopping times, including $\tau=\infty$, and interpret the discounted payoff as zero on $\{\tau=\infty\}$.

\begin{de}
For initial stock price $S_0=x$, the perpetual American put value is
\[
v_*(x)
=\sup_{\tau\in\mathcal{T}}
 \widetilde{\mathbb{E}}_x
 \left[e^{-r\tau}(K-S_\tau)^+\right].
\]
\end{de}

We first recall a first-passage formula.

\begin{thm}[Laplace transform of a drifted-Brownian first-passage time]
Let
\[
X_t=\mu t+\widetilde W_t,
\qquad
\tau_m=\inf\{t\geq0:X_t=m\},
\quad m>0.
\]
Then, for $\lambda>0$,
\[
\widetilde{\mathbb{E}}\!\left[e^{-\lambda\tau_m}\right]
=\exp\!\left[-m\left(-\mu+\sqrt{\mu^2+2\lambda}\right)\right],
\]
where $e^{-\lambda\tau_m}$ is interpreted as zero when $\tau_m=\infty$.
\end{thm}

Fix a candidate exercise threshold $L\in(0,K)$. The threshold strategy exercises immediately when $x\leq L$ and otherwise exercises at
\[
\tau_L=\inf\{t\geq0:S_t=L\}.
\]

\begin{thm}
The value of the threshold-$L$ exercise strategy is
\[
v_L(x)
=\begin{cases}
K-x, & 0\leq x\leq L,\\[4pt]
(K-L)\left(\dfrac{x}{L}\right)^{-\frac{2r}{\sigma^2}},
& x>L.
\end{cases}
\]
\end{thm}

\begin{proof}
For $x>L$,
\[
S_t=x\exp\!\left(
(r-\tfrac12\sigma^2)t+\sigma\widetilde W_t
\right).
\]
The event $S_t=L$ is equivalent to
\[
-\widetilde W_t
-\frac{r-\frac12\sigma^2}{\sigma}t
=\frac1\sigma\log\frac{x}{L}.
\]
Apply the first-passage theorem with
\[
\mu=-\frac{r-\frac12\sigma^2}{\sigma},
\qquad
m=\frac1\sigma\log\frac{x}{L},
\qquad
\lambda=r.
\]
Since
\[
-\mu+\sqrt{\mu^2+2r}=\frac{2r}{\sigma},
\]
we obtain
\[
\widetilde{\mathbb{E}}_x[e^{-r\tau_L}]
=\left(\frac{x}{L}\right)^{-\frac{2r}{\sigma^2}}.
\]
At exercise, the payoff is $K-L$, which proves the formula.
\end{proof}

For fixed $x>L$, maximizing $v_L(x)$ is equivalent to maximizing
\[
g(L)=(K-L)L^{\frac{2r}{\sigma^2}}.
\]
The first-order condition gives the optimal threshold
\[
\boxed{L_*=\frac{2r}{2r+\sigma^2}K.}
\]
Thus the candidate value is
\[
v(x)=v_{L_*}(x)
=\begin{cases}
K-x, & 0\leq x\leq L_*,\\[4pt]
(K-L_*)\left(\dfrac{x}{L_*}\right)^{-\frac{2r}{\sigma^2}},
& x>L_*.
\end{cases}
\]
It satisfies value matching and smooth pasting:
\[
v(L_*)=K-L_*,
\qquad
v'(L_*-) = v'(L_*+)=-1.
\]
Its derivatives away from $L_*$ are
\[
v'(x)
=\begin{cases}
-1, & 0<x<L_*,\\[4pt]
-\dfrac{2r(K-L_*)}{\sigma^2x}
 \left(\dfrac{x}{L_*}\right)^{-\frac{2r}{\sigma^2}},
&x>L_*,
\end{cases}
\]
and
\[
v''(x)
=\begin{cases}
0, & 0<x<L_*,\\[4pt]
\dfrac{2r(2r+\sigma^2)(K-L_*)}{\sigma^4x^2}
 \left(\dfrac{x}{L_*}\right)^{-\frac{2r}{\sigma^2}},
&x>L_*.
\end{cases}
\]
The second derivative has a jump at $L_*$, but $v$ is continuously differentiable.

Let
\[
\mathcal{L}f(x)
=rx f'(x)+\frac12\sigma^2x^2f''(x)
\]
be the risk-neutral generator. The value function satisfies the complementarity conditions
\[
\begin{aligned}
v(x)&\geq(K-x)^+,\\
rv(x)-\mathcal{L}v(x)&\geq0,\\
\bigl(v(x)-(K-x)^+\bigr)
\bigl(rv(x)-\mathcal{L}v(x)\bigr)&=0.
\end{aligned}
\]
Equivalently,
\[
\min\bigl\{v(x)-(K-x)^+,\ rv(x)-\mathcal{L}v(x)\bigr\}=0.
\]

\begin{thm}[Verification for the perpetual put]
Let
\[
\tau_*=\inf\{t\geq0:S_t\leq L_*\}.
\]
Then $e^{-rt}v(S_t)$ is a supermartingale, and
\[
e^{-r(t\wedge\tau_*)}v(S_{t\wedge\tau_*})
\]
is a martingale. Consequently,
\[
v(x)
=\sup_{\tau\in\mathcal{T}}
 \widetilde{\mathbb{E}}_x
 \left[e^{-r\tau}(K-S_\tau)^+\right],
\]
and $\tau_*$ is optimal.
\end{thm}

\begin{proof}
Because $v\in C^1$ and is piecewise $C^2$, the generalized It\^o formula applies. No local-time term appears because the first derivative is continuous at $L_*$. We obtain
\[
\begin{aligned}
d\bigl(e^{-rt}v(S_t)\bigr)
&=e^{-rt}\bigl(\mathcal{L}v(S_t)-rv(S_t)\bigr)dt\\
&\quad+e^{-rt}\sigma S_tv'(S_t)d\widetilde W_t\\
&=-e^{-rt}rK\mathbf{1}_{\{S_t<L_*\}}dt
 +e^{-rt}\sigma S_tv'(S_t)d\widetilde W_t.
\end{aligned}
\]
The nonpositive finite-variation term proves the supermartingale property. Before $\tau_*$, the stock is in the continuation region and the drift is zero, so the stopped process is a martingale.

For any stopping time $\tau$, optional sampling and $v\geq(K-\cdot)^+$ give
\[
v(x)
\geq\widetilde{\mathbb{E}}_x
 \left[e^{-r\tau}v(S_\tau)\right]
\geq\widetilde{\mathbb{E}}_x
 \left[e^{-r\tau}(K-S_\tau)^+\right],
\]
after localization and passage to the limit. Since $0\leq v\leq K$, dominated convergence is available. For $\tau=\tau_*$, the stopped martingale and value matching give equality:
\[
v(x)
=\widetilde{\mathbb{E}}_x
 \left[e^{-r\tau_*}(K-S_{\tau_*})^+\right].
\]
\end{proof}

The supermartingale decomposition also gives a hedging interpretation.

\begin{cor}
Start with capital $X_0=v(S_0)$, hold
\[
\Delta_t=v'(S_t)
\]
shares of stock, and consume at rate
\[
C_t=rK\mathbf{1}_{\{S_t<L_*\}}.
\]
Then, until exercise,
\[
X_t=v(S_t)\geq(K-S_t)^+.
\]
Thus the strategy finances a short position in the perpetual American put regardless of when the holder exercises.
\end{cor}

\begin{proof}
A self-financing strategy with consumption satisfies
\[
dX_t=\Delta_t\,dS_t+r(X_t-\Delta_tS_t)dt-C_tdt.
\]
Hence
\[
d(e^{-rt}X_t)
=e^{-rt}\Delta_t\sigma S_t\,d\widetilde W_t-e^{-rt}C_tdt.
\]
With the stated $\Delta$ and $C$, this equals the differential of $e^{-rt}v(S_t)$. Equality of initial values then implies $X_t=v(S_t)$ up to exercise.
\end{proof}

If the stock is below $L_*$ and the holder irrationally delays exercise, then $v(S_t)=K-S_t$, $\Delta_t=-1$, and the hedger has $K$ in the money-market account. The interest $rKdt$ can be consumed while preserving the hedge.

\subsection{American options: finite maturity}

Let $0\leq t\leq T$ and $S_t=x$. Denote by $\mathcal{T}_{t,T}$ the stopping times taking values in $[t,T]$.

\begin{de}
The finite-maturity American put value is
\[
v(t,x)
=\sup_{\tau\in\mathcal{T}_{t,T}}
 \widetilde{\mathbb{E}}_{t,x}
 \left[e^{-r(\tau-t)}(K-S_\tau)^+\right].
\]
\end{de}

Define the stopping and continuation regions by
\[
\mathcal{S}
=\{(t,x):v(t,x)=(K-x)^+\},
\qquad
\mathcal{C}
=\{(t,x):v(t,x)>(K-x)^+\}.
\]
Subject to the usual regularity, $v$ satisfies the variational inequality
\[
\begin{aligned}
v(t,x)&\geq(K-x)^+,\\
rv-v_t-rxv_x-\frac12\sigma^2x^2v_{xx}&\geq0,\\
\bigl(v-(K-x)^+\bigr)
\left(rv-v_t-rxv_x-\frac12\sigma^2x^2v_{xx}\right)&=0,
\end{aligned}
\]
with terminal condition
\[
v(T,x)=(K-x)^+.
\]
Equivalently,
\[
\min\!\left\{
 v-(K-x)^+,
 rv-v_t-rxv_x-\frac12\sigma^2x^2v_{xx}
\right\}=0.
\]

\begin{thm}[Verification principle]
Assume $v$ has sufficient regularity for the generalized It\^o formula, satisfies the preceding variational inequality and growth conditions, and define
\[
\tau_*
=\inf\{u\in[t,T]:(u,S_u)\in\mathcal{S}\}.
\]
Then
\[
e^{-r(u-t)}v(u,S_u),
\qquad t\leq u\leq T,
\]
is a supermartingale, and the stopped process
\[
e^{-r((u\wedge\tau_*)-t)}
 v(u\wedge\tau_*,S_{u\wedge\tau_*})
\]
is a martingale. In particular, $\tau_*$ is optimal.
\end{thm}

\begin{proof}
It\^o's formula gives
\[
\begin{aligned}
d\bigl(e^{-r(u-t)}v(u,S_u)\bigr)
&=e^{-r(u-t)}
 \left(v_u+rS_uv_x+\frac12\sigma^2S_u^2v_{xx}-rv\right)du\\
&\quad+e^{-r(u-t)}\sigma S_uv_x\,d\widetilde W_u.
\end{aligned}
\]
The variational inequality makes the drift nonpositive everywhere and zero in $\mathcal{C}$. The process stopped on first entry into $\mathcal{S}$ therefore has zero drift. Optional sampling, the payoff domination, and terminal value matching complete the verification argument.
\end{proof}

For the standard American put, the stopping region is typically of the form
\[
\mathcal{S}=\{(t,x):x\leq b(t)\}
\]
for an exercise boundary $b(t)$. If a holder delays exercise after the process enters the stopping region, the seller can continue the delta hedge
\[
\Delta_u=v_x(u,S_u)
\]
and consume at rate
\[
C_u=rK\mathbf{1}_{\{(u,S_u)\in\mathcal{S}\}}.
\]
The corresponding wealth satisfies
\[
X_u=v(u,S_u)\geq(K-S_u)^+
\]
until exercise or maturity.

\subsection{American calls on non-dividend-paying stocks}

The following convexity argument explains why a non-dividend-paying American call should not be exercised early.

\begin{lem}
Let $h:[0,\infty)\to[0,\infty)$ be convex with $h(0)=0$, and assume the relevant integrability. Then
\[
e^{-rt}h(S_t)
\]
is a submartingale under $\widetilde{\mathbb{P}}$.
\end{lem}

\begin{proof}
Convexity and $h(0)=0$ imply
\[
h(\lambda x)\leq\lambda h(x),
\qquad 0\leq\lambda\leq1.
\]
For $0\leq u\leq t$, set $\lambda=e^{-r(t-u)}$. Then conditional Jensen's inequality and the martingale property of $e^{-rt}S_t$ give
\[
\begin{aligned}
\widetilde{\mathbb{E}}\!\left[
 e^{-r(t-u)}h(S_t)\mid\mathcal{F}_u
\right]
&\geq
\widetilde{\mathbb{E}}\!\left[
 h(e^{-r(t-u)}S_t)\mid\mathcal{F}_u
\right]\\
&\geq
h\!\left(
 \widetilde{\mathbb{E}}[e^{-r(t-u)}S_t\mid\mathcal{F}_u]
\right)\\
&=h(S_u).
\end{aligned}
\]
Multiplying by $e^{-ru}$ yields the submartingale property.
\end{proof}

\begin{thm}
Let $h$ satisfy the assumptions of the lemma. The American claim with exercise payoff $h(S_t)$ and maturity $T$ has the same value as the European claim paying $h(S_T)$:
\[
\sup_{\tau\in\mathcal{T}_{t,T}}
\widetilde{\mathbb{E}}\!\left[
 e^{-r(\tau-t)}h(S_\tau)\mid\mathcal{F}_t
\right]
=
\widetilde{\mathbb{E}}\!\left[
 e^{-r(T-t)}h(S_T)\mid\mathcal{F}_t
\right].
\]
In particular, an American call on a non-dividend-paying stock should not be exercised before maturity.
\end{thm}

\begin{proof}
For every stopping time $\tau\in\mathcal{T}_{t,T}$, optional sampling for the submartingale $e^{-ru}h(S_u)$ gives
\[
\widetilde{\mathbb{E}}\!\left[
 e^{-r\tau}h(S_\tau)\mid\mathcal{F}_t
\right]
\leq
\widetilde{\mathbb{E}}\!\left[
 e^{-rT}h(S_T)\mid\mathcal{F}_t
\right].
\]
Equality is achieved by choosing $\tau=T$. The call payoff $h(x)=(x-K)^+$ is nonnegative, convex, and satisfies $h(0)=0$.
\end{proof}

    \section{Term-structure models}

\subsection{Introduction}

Fixed-income markets quote bond prices rather than a single directly observable interest rate. A coupon-bearing bond can be decomposed into zero-coupon cash flows, and the prices of those zero-coupon bonds determine yields and forward rates.

Let $P(t,T)$ denote the time-$t$ price of a zero-coupon bond that pays one unit of currency at maturity $T$, where $0\leq t\leq T$. Its continuously compounded yield to maturity is
\[
y(t,T)
=-\frac{\log P(t,T)}{T-t},
\qquad t<T,
\]
so that
\[
P(t,T)=e^{-y(t,T)(T-t)}.
\]

If a bond pays deterministic cash flows $C_i$ at dates
\[
t<T_1<\cdots<T_n,
\]
then its no-arbitrage price is
\[
\Pi_t=\sum_{i=1}^n C_iP(t,T_i).
\]
Observed coupon-bond prices can therefore be used to bootstrap zero-coupon prices. For example, if bond $j$ pays $C_{j,i}$ at $T_i$, with $C_{j,j}\neq0$, then the triangular system
\[
\Pi_j=\sum_{i=1}^j C_{j,i}P(t,T_i),
\qquad j=1,\ldots,n,
\]
can be solved successively for $P(t,T_1),\ldots,P(t,T_n)$.

The map
\[
T\longmapsto y(t,T)
\]
is the yield curve observed at time $t$. It may be upward sloping, downward sloping, flat, or humped; none of these shapes is universal.

For $t\leq T_1<T_2$, the simple forward rate fixed at time $t$ for lending over $[T_1,T_2]$ is
\[
L(t;T_1,T_2)
=\frac{1}{T_2-T_1}
 \left(\frac{P(t,T_1)}{P(t,T_2)}-1\right).
\]
Indeed, buying one $T_1$-bond and financing it by shorting
$P(t,T_1)/P(t,T_2)$ units of the $T_2$-bond creates a loan from $T_1$ to $T_2$.

The instantaneous forward rate is
\[
f(t,T)
=-\frac{\partial}{\partial T}\log P(t,T),
\]
which implies
\[
P(t,T)
=\exp\!\left(-\int_t^Tf(t,u)\,du\right).
\]
The short rate is
\[
r_t=f(t,t).
\]

There are two principal modeling approaches:
\begin{enumerate}
    \item specify the dynamics of the short rate or of a finite collection of state variables and derive bond prices; or
    \item specify the dynamics of the entire forward-rate curve, as in the Heath--Jarrow--Morton framework.
\end{enumerate}

\subsection{Affine term-structure models}

\subsubsection*{One-factor models}

The Vasicek model specifies, under a risk-neutral measure $\widetilde{\mathbb{P}}$,
\[
dr_t=\kappa(\theta-r_t)dt+\sigma\,d\widetilde W_t,
\qquad
\kappa,\theta,\sigma>0.
\]
This Gaussian model is mean reverting but permits negative rates. The time-inhomogeneous extension
\[
dr_t=\bigl(\theta(t)-\kappa r_t\bigr)dt+\sigma\,d\widetilde W_t
\]
is commonly called the one-factor Hull--White model.

The Cox--Ingersoll--Ross model is
\[
dr_t=\kappa(\theta-r_t)dt+\sigma\sqrt{r_t}\,d\widetilde W_t,
\qquad
\kappa,\theta,\sigma>0.
\]
Starting from $r_0\geq0$, the process remains nonnegative. The Feller condition
\[
2\kappa\theta\geq\sigma^2
\]
makes the boundary at zero unattainable when $r_0>0$, and hence ensures strict positivity.

Suppose generally that, under $\widetilde{\mathbb{P}}$,
\[
dr_t=a(t,r_t)dt+b(t,r_t)d\widetilde W_t.
\]
The zero-coupon bond price is
\[
P(t,T)
=\widetilde{\mathbb{E}}\!\left[
 \exp\!\left(-\int_t^Tr_sds\right)
 \bigg|\mathcal{F}_t
\right]
=v(t,r_t;T).
\]
By Feynman--Kac, $v$ solves
\[
\begin{cases}
 v_t+a(t,r)v_r+\dfrac12b(t,r)^2v_{rr}-rv=0,
    & t<T,\\[4pt]
 v(T,r;T)=1.
\end{cases}
\]
In affine models, the solution has exponential-affine form
\[
P(t,T)=A(t,T)e^{-C(t,T)r_t},
\]
so the yield
\[
y(t,T)
=-\frac{\log A(t,T)}{T-t}
 +\frac{C(t,T)}{T-t}r_t
\]
is affine in the current short rate.

\subsubsection*{Two-factor Gaussian models}

A two-factor Vasicek model may be written as
\[
\begin{aligned}
dX_1(t)
&=\bigl(a_1-b_{11}X_1(t)-b_{12}X_2(t)\bigr)dt
 +\sigma_1d\widetilde W_1(t),\\
dX_2(t)
&=\bigl(a_2-b_{21}X_1(t)-b_{22}X_2(t)\bigr)dt
 +\sigma_2d\widetilde W_2(t),\\
r_t&=\delta_0+\delta_1X_1(t)+\delta_2X_2(t),
\end{aligned}
\]
where
\[
d\widetilde W_1(t)d\widetilde W_2(t)=\rho\,dt,
\qquad -1<\rho<1.
\]
After a linear transformation, one may often use a canonical form driven by independent Brownian motions, such as
\[
\begin{aligned}
dY_1(t)&=-\lambda_1Y_1(t)dt+d\widetilde B_1(t),\\
dY_2(t)&=-\lambda_{21}Y_1(t)dt-\lambda_2Y_2(t)dt+d\widetilde B_2(t),\\
r_t&=\delta_0+\delta_1Y_1(t)+\delta_2Y_2(t).
\end{aligned}
\]

\subsubsection*{Two-factor square-root models}

A canonical two-factor square-root model is
\[
\begin{aligned}
dY_1(t)
&=\bigl(\mu_1-\lambda_{11}Y_1(t)-\lambda_{12}Y_2(t)\bigr)dt
 +\sqrt{Y_1(t)}\,d\widetilde W_1(t),\\
dY_2(t)
&=\bigl(\mu_2-\lambda_{21}Y_1(t)-\lambda_{22}Y_2(t)\bigr)dt
 +\sqrt{Y_2(t)}\,d\widetilde W_2(t),\\
r_t&=\delta_0+\delta_1Y_1(t)+\delta_2Y_2(t).
\end{aligned}
\]
Conditions such as
\[
\mu_1,\mu_2\geq0,
\qquad
\lambda_{12},\lambda_{21}\leq0
\]
make the drift point inward on the boundary of the nonnegative quadrant and help preserve nonnegativity. Mean reversion is a separate stability property: it requires the drift matrix
\[
\Lambda=
\begin{pmatrix}
\lambda_{11}&\lambda_{12}\\
\lambda_{21}&\lambda_{22}
\end{pmatrix}
\]
to have eigenvalues with positive real parts. Additional Feller-type restrictions are used when strict positivity is required.

A mixed square-root/Gaussian specification is
\[
\begin{aligned}
dY_1(t)
&=\bigl(\mu-\lambda_1Y_1(t)\bigr)dt
 +\sqrt{Y_1(t)}\,d\widetilde W_1(t),\\
dY_2(t)
&=-\lambda_2Y_2(t)dt
 +\sigma_{21}\sqrt{Y_1(t)}\,d\widetilde W_1(t)
 +\sqrt{\alpha+\beta Y_1(t)}\,d\widetilde W_2(t),\\
r_t&=\delta_0+\delta_1Y_1(t)+\delta_2Y_2(t),
\end{aligned}
\]
where $\widetilde W_1$ and $\widetilde W_2$ are independent and the parameters are chosen so that the square roots are well defined.

In affine multi-factor models, the bond price generally has the form
\[
P(t,T)
=\exp\!\left(
 A(t,T)-C_1(t,T)Y_1(t)-C_2(t,T)Y_2(t)
\right),
\]
with deterministic coefficient functions obtained from Riccati equations. The yield is affine in the factors.

\subsection{Heath--Jarrow--Morton models}

The Heath--Jarrow--Morton approach models the full instantaneous-forward curve. Under the physical measure $\mathbb{P}$, suppose
\[
df(t,T)
=\alpha(t,T)dt+\sigma(t,T)^\top dW_t,
\qquad 0\leq t\leq T,
\]
where $W$ is $d$-dimensional and $\sigma(t,T)\in\mathbb{R}^d$.

Define
\[
\alpha^*(t,T)=\int_t^T\alpha(t,u)\,du,
\qquad
\sigma^*(t,T)=\int_t^T\sigma(t,u)\,du.
\]
Let
\[
g(t,T)=\int_t^Tf(t,u)\,du,
\qquad
P(t,T)=e^{-g(t,T)}.
\]
Leibniz's rule gives
\[
dg(t,T)
=\bigl(-r_t+\alpha^*(t,T)\bigr)dt
 +\sigma^*(t,T)^\top dW_t.
\]
It\^o's formula therefore yields
\[
\boxed{
\frac{dP(t,T)}{P(t,T)}
=\left(
 r_t-\alpha^*(t,T)
 +\frac12\|\sigma^*(t,T)\|^2
 \right)dt
-\sigma^*(t,T)^\top dW_t.}
\]

Let $\theta_t\in\mathbb{R}^d$ be the market price of Brownian risk and define
\[
dW_t^{\widetilde{\mathbb{P}}}=dW_t+\theta_tdt.
\]
Under the corresponding equivalent measure, substitution gives
\[
\frac{dP(t,T)}{P(t,T)}
=\left(
 r_t-\alpha^*(t,T)
 +\frac12\|\sigma^*(t,T)\|^2
 +\sigma^*(t,T)^\top\theta_t
 \right)dt
-\sigma^*(t,T)^\top dW_t^{\widetilde{\mathbb{P}}}.
\]
For discounted bond prices to be local martingales, the HJM drift restriction under $\mathbb{P}$ is
\[
\boxed{
\alpha^*(t,T)
=\sigma^*(t,T)^\top\theta_t
 +\frac12\|\sigma^*(t,T)\|^2.}
\]
Differentiating with respect to $T$ gives
\[
\alpha(t,T)
=\sigma(t,T)^\top\theta_t
 +\sigma(t,T)^\top\sigma^*(t,T).
\]
Hence, under the risk-neutral measure,
\[
\boxed{
df(t,T)
=\sigma(t,T)^\top\sigma^*(t,T)dt
 +\sigma(t,T)^\top dW_t^{\widetilde{\mathbb{P}}}.}
\]
The corresponding bond dynamics are
\[
\boxed{
\frac{dP(t,T)}{P(t,T)}
=r_tdt-\sigma^*(t,T)^\top dW_t^{\widetilde{\mathbb{P}}}.}
\]
Thus, once the forward-rate volatility is specified, absence of arbitrage determines the risk-neutral drift.

\section{Jump models}

\subsection{Review of Poisson processes}

Diffusion models have continuous paths and therefore cannot generate sudden price moves. Jump models can capture discontinuities, heavier tails, and some features of implied-volatility smiles and skews.

Let $\xi_1,\xi_2,\ldots$ be independent exponential random variables with parameter $\lambda>0$:
\[
\mathbb{P}(\xi_i>t)=e^{-\lambda t},
\qquad
\mathbb{E}[\xi_i]=\frac1\lambda.
\]
The exponential distribution is memoryless:
\[
\mathbb{P}(\xi_i>s+t\mid\xi_i>s)
=\mathbb{P}(\xi_i>t).
\]
Define the $n$th arrival time by
\[
T_n=\sum_{i=1}^n\xi_i
\]
and the counting process by
\[
N_t=\max\{n\geq0:T_n\leq t\}.
\]
Then $N$ is a Poisson process of intensity $\lambda$ and
\[
\mathbb{P}(N_t=k)
=e^{-\lambda t}\frac{(\lambda t)^k}{k!},
\qquad
\mathbb{E}[N_t]=\operatorname{Var}(N_t)=\lambda t.
\]
The arrival time $T_n$ has the Erlang density
\[
f_{T_n}(s)
=\frac{\lambda^n s^{n-1}}{(n-1)!}e^{-\lambda s},
\qquad s>0.
\]
The compensated process
\[
M_t=N_t-\lambda t
\]
is a martingale.

Let $Y_1,Y_2,\ldots$ be i.i.d., independent of $N$, and suppose $\mathbb{E}|Y_1|<\infty$. The compound Poisson process is
\[
Q_t=\sum_{i=1}^{N_t}Y_i.
\]
A standard compound Poisson process has at most one arrival at any given time almost surely; the jump at an arrival may have a random size $Y_i$. If
\[
\beta=\mathbb{E}[Y_1],
\]
then
\[
Q_t-\lambda\beta t
\]
is a martingale. If $M_Y(u)=\mathbb{E}[e^{uY_1}]$ is finite, then
\[
\mathbb{E}[e^{uQ_t}]
=\exp\!\left(\lambda t\bigl(M_Y(u)-1\bigr)\right).
\]

\subsection{Stochastic calculus for jump processes}

For a compound Poisson process with jump sizes $Y_i$, define the random jump measure by
\[
J((0,t]\times A)
=\sum_{i=1}^{N_t}\mathbf{1}_{\{Y_i\in A\}},
\]
for Borel sets $A\subseteq\mathbb{R}$. Then
\[
Q_t=\int_0^t\int_{\mathbb{R}}y\,J(ds,dy).
\]
If the jump-size distribution is $G(dy)$, the compensator is
\[
\nu(dy)ds=\lambda G(dy)ds,
\]
and the compensated jump measure is
\[
\widetilde J(ds,dy)=J(ds,dy)-\lambda G(dy)ds.
\]
For predictable integrands satisfying the appropriate integrability condition,
\[
\int_0^t\int H_s(y)\,\widetilde J(ds,dy)
\]
is a local martingale.

A finite-activity jump diffusion may be written as
\[
X_t
=X_0+\int_0^tb_sds+\int_0^t\gamma_s^\top dW_s
 +\int_0^t\int_{\mathbb{R}}y\,J(ds,dy).
\]
Its continuous part is
\[
X_t^c
=X_0+\int_0^tb_sds+\int_0^t\gamma_s^\top dW_s,
\]
and its jump at time $s$ is
\[
\Delta X_s=X_s-X_{s-}.
\]
Stochastic integrands are evaluated predictably, using left limits at jumps. For example,
\[
\int_0^tH_s\,dX_s
=\int_0^tH_s b_sds
 +\int_0^tH_s\gamma_s^\top dW_s
 +\sum_{0<s\leq t}H_s\Delta X_s,
\]
where $H$ is predictable. If $X$ is a local martingale and $H$ is $X$-integrable, then $H\mathbin{\boldsymbol\cdot}X$ is a local martingale; predictability alone does not make an integral with respect to an arbitrary semimartingale a martingale.

For semimartingales $X$ and $Y$,
\[
[X,Y]_t
=[X^c,Y^c]_t
 +\sum_{0<s\leq t}\Delta X_s\Delta Y_s.
\]
In particular,
\[
[X]_t=[X^c]_t+\sum_{0<s\leq t}(\Delta X_s)^2.
\]

\begin{thm}[It\^o formula with jumps]
Let $f\in C^2(\mathbb{R})$ and let $X$ be a semimartingale. Then
\[
\begin{aligned}
f(X_t)
&=f(X_0)
 +\int_0^tf'(X_{s-})\,dX_s
 +\frac12\int_0^tf''(X_{s-})\,d[X^c]_s\\
&\quad+\sum_{0<s\leq t}
 \left(
 f(X_s)-f(X_{s-})-f'(X_{s-})\Delta X_s
 \right).
\end{aligned}
\]
\end{thm}

For a $C^{1,2}$ function $f(t,x)$, the time-dependent form is
\[
\begin{aligned}
f(t,X_t)
&=f(0,X_0)
 +\int_0^tf_t(s,X_{s-})ds
 +\int_0^tf_x(s,X_{s-})dX_s\\
&\quad+\frac12\int_0^tf_{xx}(s,X_{s-})d[X^c]_s\\
&\quad+\sum_{0<s\leq t}
 \left(
 f(s,X_s)-f(s,X_{s-})
 -f_x(s,X_{s-})\Delta X_s
 \right).
\end{aligned}
\]

The product rule is
\[
d(X_tY_t)
=X_{t-}\,dY_t+Y_{t-}\,dX_t+d[X,Y]_t.
\]

If $X_0=0$ and $\Delta X_t>-1$, the Dol\'eans--Dade exponential is
\[
\mathcal{E}(X)_t
=\exp\!\left(X_t-\frac12[X^c]_t\right)
 \prod_{0<s\leq t}(1+\Delta X_s)e^{-\Delta X_s}.
\]
It is the unique solution of
\[
dZ_t=Z_{t-}\,dX_t,
\qquad Z_0=1.
\]
When $X$ is a local martingale, $\mathcal{E}(X)$ is a local martingale, but additional integrability is required for it to be a true martingale.

\subsection{Changes of measure for Poisson processes}

Suppose $N$ is Poisson with intensity $\lambda$ under $\mathbb{P}$, and let $\widetilde\lambda>0$. The process
\[
Z_t
=\exp\!\bigl((\lambda-\widetilde\lambda)t\bigr)
 \left(\frac{\widetilde\lambda}{\lambda}\right)^{N_t}
\]
is a martingale. Under the measure defined by
\[
\frac{d\widetilde{\mathbb{P}}}{d\mathbb{P}}
\bigg|_{\mathcal{F}_T}=Z_T,
\]
the process $N$ is Poisson with intensity $\widetilde\lambda$.

Now suppose the marks have density $g$ under $\mathbb{P}$, and choose a target density $\widetilde g$ that is absolutely continuous with respect to $g$; equivalently,
\[
g(y)=0\implies\widetilde g(y)=0
\quad\text{for almost every }y.
\]
Then
\[
Z_t
=\exp\!\bigl((\lambda-\widetilde\lambda)t\bigr)
 \prod_{i=1}^{N_t}
 \frac{\widetilde\lambda\,\widetilde g(Y_i)}
      {\lambda g(Y_i)}
\]
changes the intensity to $\widetilde\lambda$ and the mark density to $\widetilde g$, provided $Z$ is a true martingale. The new measure is equivalent to $\mathbb{P}$ precisely when $g$ and $\widetilde g$ have the same support up to null sets.

A Brownian change of drift can be combined with this jump change of measure. If
\[
Z_t^{W}
=\exp\!\left(
 -\int_0^t\theta_s^\top dW_s
 -\frac12\int_0^t\|\theta_s\|^2ds
\right)
\]
and $Z^J$ is the jump likelihood above, then, under appropriate martingale and dependence assumptions, $Z_t=Z_t^WZ_t^J$ defines a measure under which
\[
\widetilde W_t=W_t+\int_0^t\theta_sds
\]
is Brownian and the jump process has intensity $\widetilde\lambda$ and mark density $\widetilde g$.

\subsection{Jump-diffusion asset models}

Let
\[
Q_t=\sum_{i=1}^{N_t}Y_i,
\qquad Y_i>-1,
\qquad
\beta=\mathbb{E}[Y_i].
\]
Under the physical measure, consider
\[
\frac{dS_t}{S_{t-}}
=\alpha\,dt+\sigma\,dW_t+d(Q_t-\lambda\beta t).
\]
Equivalently,
\[
\frac{dS_t}{S_{t-}}
=(\alpha-\lambda\beta)dt+\sigma dW_t+dQ_t.
\]
Its solution is
\[
S_t
=S_0\exp\!\left(
 (\alpha-\lambda\beta-\tfrac12\sigma^2)t
 +\sigma W_t
\right)
\prod_{i=1}^{N_t}(1+Y_i).
\]

Under an equivalent measure $\widetilde{\mathbb{P}}$, suppose the Brownian kernel is $\theta$, the jump intensity is $\widetilde\lambda$, the mark density is $\widetilde g$, and
\[
\widetilde\beta
=\widetilde{\mathbb{E}}[Y].
\]
With the convention
\[
d\widetilde W_t=dW_t+\theta_tdt,
\]
the discounted stock is a local martingale exactly when
\[
\boxed{
\alpha-\lambda\beta-\sigma\theta_t
=r-\widetilde\lambda\widetilde\beta.}
\]
Equivalently, under $\widetilde{\mathbb{P}}$,
\[
\frac{dS_t}{S_{t-}}
=r\,dt+\sigma\,d\widetilde W_t
 +\int_{(-1,\infty)}y\,\widetilde J(dt,dy),
\]
where
\[
\widetilde J(dt,dy)
=J(dt,dy)-\widetilde\lambda\widetilde g(y)dy\,dt.
\]
In raw-jump form,
\[
\frac{dS_t}{S_{t-}}
=(r-\widetilde\lambda\widetilde\beta)dt
 +\sigma d\widetilde W_t
 +\int_{(-1,\infty)}y\,J(dt,dy).
\]

There are generally many Brownian kernels, jump intensities, and mark densities satisfying the martingale condition. This nonuniqueness is a manifestation of market incompleteness: the stock and money-market account do not usually span all Brownian and jump risks.

\subsection{Pricing derivatives driven by jump processes}

Let $\tau=T-t$. Under $\widetilde{\mathbb{P}}$,
\[
S_T
=S_t\exp\!\left(
(r-\widetilde\lambda\widetilde\beta-\tfrac12\sigma^2)\tau
+\sigma(\widetilde W_T-\widetilde W_t)
\right)
\prod_{i=N_t+1}^{N_T}(1+Y_i).
\]
Define, conditional on $N_T-N_t=j$,
\[
R_j
=e^{-\widetilde\lambda\widetilde\beta\tau}
 \prod_{i=1}^j(1+Y_i),
\qquad R_0=e^{-\widetilde\lambda\widetilde\beta\tau}.
\]
Conditional on the number and sizes of jumps, the remaining uncertainty is Black--Scholes lognormal risk. Therefore the European call price is the Poisson mixture
\[
\boxed{
\begin{aligned}
c(t,S_t)
&=\sum_{j=0}^{\infty}
 e^{-\widetilde\lambda\tau}
 \frac{(\widetilde\lambda\tau)^j}{j!}
 \widetilde{\mathbb{E}}\!\left[
 C_{\mathrm{BS}}(\tau,S_tR_j;K,r,\sigma)
 \right],
\end{aligned}}
\]
where the expectation inside the sum is over $j$ independent marks with density $\widetilde g$. The notation $C_{\mathrm{BS}}$ denotes the Black--Scholes call-pricing function; its arguments are time to maturity $\tau$, spot $s$, strike $K$, rate $r$, and volatility $\sigma$.

\subsubsection*{Partial integro-differential equation}

Assume $c(t,x)$ is sufficiently smooth and has terminal condition
\[
c(T,x)=(x-K)^+.
\]
The risk-neutral generator of the stock in raw-jump form is
\[
\begin{aligned}
\mathcal{L}c(t,x)
&=(r-\widetilde\lambda\widetilde\beta)x c_x(t,x)
 +\frac12\sigma^2x^2c_{xx}(t,x)\\
&\quad+\widetilde\lambda
 \int_{-1}^{\infty}
 \bigl[c(t,x(1+y))-c(t,x)\bigr]
 \widetilde g(y)\,dy.
\end{aligned}
\]
The discounted option value is a martingale, so $c$ solves
\[
\boxed{
\begin{aligned}
0={}&c_t(t,x)
 +(r-\widetilde\lambda\widetilde\beta)x c_x(t,x)
 +\frac12\sigma^2x^2c_{xx}(t,x)-rc(t,x)\\
&+\widetilde\lambda
 \int_{-1}^{\infty}
 \bigl[c(t,x(1+y))-c(t,x)\bigr]
 \widetilde g(y)\,dy,
\end{aligned}}
\]
with $c(T,x)=(x-K)^+$.

Equivalently, using the compensated jump measure, the PIDE may be written
\[
\begin{aligned}
0={}&c_t+rx c_x+\frac12\sigma^2x^2c_{xx}-rc\\
&+\widetilde\lambda
 \int_{-1}^{\infty}
 \bigl[c(t,x(1+y))-c(t,x)-xyc_x(t,x)\bigr]
 \widetilde g(y)\,dy.
\end{aligned}
\]

\subsubsection*{Delta hedging and incompleteness}

Suppose a self-financing portfolio holds $\Gamma_t$ shares of stock. In discounted terms,
\[
\begin{aligned}
d(e^{-rt}X_t)
&=e^{-rt}\Gamma_t\sigma S_{t-}\,d\widetilde W_t\\
&\quad+e^{-rt}\Gamma_tS_{t-}
 \int_{-1}^{\infty}y\,\widetilde J(dt,dy).
\end{aligned}
\]
The discounted option value satisfies
\[
\begin{aligned}
d(e^{-rt}c(t,S_t))
&=e^{-rt}\sigma S_{t-}c_x(t,S_{t-})\,d\widetilde W_t\\
&\quad+e^{-rt}
 \int_{-1}^{\infty}
 \bigl[c(t,S_{t-}(1+y))-c(t,S_{t-})\bigr]
 \widetilde J(dt,dy).
\end{aligned}
\]
Choosing
\[
\Gamma_t=c_x(t,S_{t-})
\]
removes the Brownian mismatch, but the jump mismatch remains:
\[
\begin{aligned}
d\bigl(e^{-rt}(X_t-c(t,S_t))\bigr)
=e^{-rt}\int_{-1}^{\infty}H_t(y)\,\widetilde J(dt,dy),
\end{aligned}
\]
where
\[
H_t(y)
=S_{t-}y\,c_x(t,S_{t-})
-\bigl[c(t,S_{t-}(1+y))-c(t,S_{t-})\bigr].
\]
For a convex option value,
\[
c(t,x(1+y))-c(t,x)
\geq xy\,c_x(t,x),
\]
so
\[
H_t(y)\leq0.
\]
Consequently, between jumps the compensator makes the delta-hedging error drift upward, so the stock hedge tends to outperform the option. At a jump, the error changes by $H_t(Y)\leq0$, so the hedge underperforms. A single stock position cannot in general eliminate both Brownian and jump-size risk; this is another direct demonstration that the jump-diffusion market is incomplete.

\end{document}